\chapter{接触力学与接触约束优化——摩擦锥 / GIWC / 互补约束 / 质心动量}
\label{ch:contact-mechanics}

% ════════════════════════════════════════════════════════════════
%  控制部 C3 · 接触力学深化章。卷五《机器人最优控制》腿足/接触子部。
%  源：Tutorial 第 52 章「接触力学与约束优化」(2048 行)。
%  与 ch:control 的硬边界：本章是其第 6--7 节(WBC/Cheetah)的理论深化，
%  只补误差界/wrench cone/GIWC/LCP-Delassus/三路线/CCRBA/角动量可控性，
%  不重述 WBC-QP 搭建与 Cheetah 凸 MPC 概览。
%  记号铁律：wrench 用本书序 w=[f;n](力在前)；外部文献(Caron/Featherstone)
%  用 [n;f]，对照处显式标注差一个 6x6 置换。
% ════════════════════════════════════════════════════════════════

\begin{practice}[前置自测]
开始之前，先试答下列问题。若已能流畅作答，可只速读本章的小结表与定理速查；若卡住，正是本章要为你解决的。
\begin{enumerate}
  \item 一个空间力（wrench）$\boldsymbol{w}\in\mathbb{R}^6$ 由力 $\boldsymbol{f}$ 与力矩 $\boldsymbol{n}$ 组成，二者如何排序才与空间速度 $\boldsymbol{\varpi}=[\boldsymbol{\nu};\boldsymbol{\omega}]$（线速度在前）配对，使功率 $\boldsymbol{w}^\top\boldsymbol{\varpi}$ 不依赖排序约定？（\cref{ch:rigid_body}）
  \item \cref{ch:control} 的全身控制把摩擦锥「线性化为方形金字塔」。这个近似是偏保守还是偏冒进？误差有多大？为什么 $k=4$ 轴对齐时会退化成一个过松的盒形锥？
  \item 刚体动力学 $\mathbf{M}\dot{\boldsymbol{\nu}}+\boldsymbol{h}=\mathbf{S}^\top\boldsymbol{\tau}+\mathbf{J}_c^\top\boldsymbol{\lambda}$ 与单边接触条件 $0\le\phi\perp\lambda^n\ge0$ 怎样耦合成一个线性互补问题（LCP）？其中那个稠密的 $\mathbf{A}=\mathbf{J}_n\mathbf{M}^{-1}\mathbf{J}_n^\top$ 矩阵是什么？
  \item 质心动量 $\boldsymbol{h}_G=[\boldsymbol{l}_G;\boldsymbol{k}_G]$（线动量在前，与本书 wrench/动量序一致）关于关节速度 $\dot{\boldsymbol{q}}$ 是线性的吗？若是，那个线性映射（质心动量矩阵 CMM）能否在 $O(N)$ 时间内构造？
  \item 一只完全离地的机器人（飞行相位）能否在空中改变自身偏航朝向？猫从倒置自由落体却能翻身平稳落地，这与角动量守恒矛盾吗？
\end{enumerate}
\end{practice}

\section*{本章目标}
学完本章，你应当能够：
\begin{enumerate}
  \item \textbf{陈述并推导}接触力的三大物理铁律——法向非负（Signorini）、库仑摩擦锥、互补性——及其变分不等式（法锥）观点；
  \item \textbf{完成}摩擦锥线性化的内/外切推导，给出面积误差界，并说明「内切才保守」这一安全性辨析；
  \item \textbf{掌握}多面体锥的 $H$-表示与 $V$-表示对偶（Minkowski--Weyl、双重描述法），并用 GIWC 给出「机器人能否静止站稳」的几何 LP 判据；
  \item \textbf{推导}从刚体动力学到 LCP 的完整链条，写出 Delassus 算子，证明其半正定，并说明可解性需附「可行」限定；
  \item \textbf{区分}互补约束的三大处理路线（预定义模式 / 软化 / 接触隐式轨迹优化）及各自的数学基础（含 Fischer--Burmeister 函数的零点等价性）；
  \item \textbf{解释}接触 Jacobian 的对偶关系 $\boldsymbol{\tau}=\mathbf{J}_c^\top\boldsymbol{\lambda}$、漂移项与 Baumgarte 稳定化；
  \item \textbf{复述} CCRBA 算法的 $O(N)$ 递推结构，理解 CMM、CCRBI 与质量矩阵的区别；
  \item \textbf{说明}角动量在飞行相位的守恒与可控性，并以猫翻身悖论辨析质心模型的盲区；
  \item \textbf{权衡}刚体接触与软接触的建模取舍，理解可微接触前沿的核心矛盾。
\end{enumerate}

\subsection*{承上：从「控制导论」到接触力学的理论地基}
\cref{ch:control} 在其全身控制（WBC）与足式凸 MPC 两节里，已经写下浮动基动力学 $\mathbf{M}\dot{\boldsymbol{\nu}}+\boldsymbol{h}=\mathbf{S}^\top\boldsymbol{\tau}+\mathbf{J}_c^\top\boldsymbol{\lambda}$、刚性接触约束 $\mathbf{J}_c\dot{\boldsymbol{\nu}}=-\dot{\mathbf{J}}_c\boldsymbol{\nu}$、把摩擦锥「线性化为金字塔」的一句话，以及 MIT Cheetah\,3 的凸 MPC 旗舰例（\cref{eq:ctrl-friction-pyramid} 即那条金字塔不等式）。那一章把这些当成\emph{已知工具}直接搭进 QP；\textbf{本章则是它的理论深化}——专门补齐那里一笔带过的深层结构：摩擦锥线性化的误差界、面接触的 wrench cone、广义惯性 wrench 锥（GIWC）与静平衡 LP 判据、刚体动力学到 LCP 的完整推导与 Delassus 算子、互补约束的三条工程路线、质心动量矩阵的 $O(N)$ 算法、以及角动量的守恒与可控性。凡涉及 WBC-QP 的搭建、刚性接触约束式、Cheetah 凸 MPC 的概览，本章一律不重述，只在用到处 \cref{ch:control}。

\subsection*{本章知识导航}
全章围绕一条主线：\emph{接触力受什么物理约束}（三大铁律）$\to$ \emph{这些约束如何进入优化器}（线性化、互补三路线、GIWC）$\to$ \emph{接触力如何在动力学里传递}（接触 Jacobian、质心动量矩阵）。

\begin{center}
\small
\begin{tikzpicture}[
  node distance=5mm and 7mm, >=Stealth,
  phy/.style={draw, rounded corners, align=center, font=\small, inner sep=3pt, fill=main!8, draw=main!60!black},
  mdl/.style={draw, rounded corners, align=center, font=\small, inner sep=3pt, fill=second!8, draw=second!60!black},
  brg/.style={draw, rounded corners, align=center, font=\small, inner sep=3pt, fill=third!8, draw=third!60!black},
  ln/.style={->, thick, gray!60!black}]
\node[phy] (law) {三大铁律\\ 法向 / 摩擦锥 / 互补};
\node[mdl, above right=2mm and 14mm of law] (lin) {线性化\\ SOC $\to$ QP};
\node[mdl, right=14mm of law] (cone) {锥的两种表示\\ GIWC 静平衡};
\node[mdl, below right=2mm and 14mm of law] (cmp) {互补三路线\\ LCP / 软化 / CITO};
\node[brg, right=12mm of cone] (jac) {接触 Jacobian\\ $\boldsymbol{\tau}=\mathbf{J}_c^\top\boldsymbol{\lambda}$};
\node[brg, below=5mm of jac] (cmm) {质心动量\\ CMM / CCRBA};
\draw[ln] (law) -- (lin); \draw[ln] (law) -- (cone); \draw[ln] (law) -- (cmp);
\draw[ln] (cone) -- (jac); \draw[ln] (cmp.east) -- (cmm.west);
\draw[ln] (jac) -- (cmm);
\end{tikzpicture}
\end{center}

\noindent\textbf{推荐路径}：\cref{sec:contact-laws}（三大铁律）是全章地基；\cref{sec:contact-linearize}（线性化）与 \cref{sec:contact-complementarity-lcp,sec:contact-routes}（互补三路线）是两条并列的「建模」支线，二者由 \cref{sec:contact-complementarity-lcp} 的 LCP 串起；\cref{sec:contact-jacobian,sec:contact-cmm} 是把接触力接回动力学的两座桥；\cref{sec:contact-angular} 是角动量可控性；\cref{sec:contact-soft} 给软/刚体接触的工程权衡。\cref{sec:contact-cone-giwc}（GIWC 静平衡判据）与 \cref{sec:contact-angular}（角动量可控性）共享同一个几何工具——\emph{接触力锥张成的可达 wrench 集合}。

\begin{note}
\textbf{本章记号与约定（务必先读）.}\;
接触力记 $\boldsymbol{\lambda}_c$，法向/切向分量 $\lambda_c^n,\boldsymbol{\lambda}_c^t$（与 \cref{ch:control} 一致，沿用不另立新符号）；接触 Jacobian $\mathbf{J}_c$，选择矩阵 $\mathbf{S}=[\mathbf{0},\mathbf{I}]$；有符号距离（gap）记 $\phi$（正为离地、零为接触、负为穿透），与重力 $\boldsymbol{g}$、Hamiltonian 区分；摩擦系数 $\mu$，线性化有效系数 $\mu_{\mathrm{eff}}$。

\textbf{空间力（wrench）排序——本章最易错处.}\;
本书空间速度取线速度在前 $\boldsymbol{\varpi}=[\boldsymbol{\nu};\boldsymbol{\omega}]$、$\se(3)$ 取平移在前 $\boldsymbol{\xi}=[\boldsymbol{\rho};\boldsymbol{\phi}]$（\cref{ch:lie,ch:rigid_body}）。与之对偶的 wrench \textbf{统一取力在前、力矩在后} $\boldsymbol{w}=[\boldsymbol{f};\boldsymbol{n}]$，使功率配对 $\boldsymbol{w}^\top\boldsymbol{\varpi}=\boldsymbol{f}^\top\boldsymbol{\nu}+\boldsymbol{n}^\top\boldsymbol{\omega}$ 与排序一致。\textbf{许多接触文献（Caron 2015、Featherstone）用力矩在前 $[\boldsymbol{n};\boldsymbol{f}]$（moment-first）}；凡与其对照，正文显式注明「外部用 $[\boldsymbol{n};\boldsymbol{f}]$、本书用 $[\boldsymbol{f};\boldsymbol{n}]$，差一个 $6\times6$ 行置换 $\mathbf{P}$」，并据此翻转其约束矩阵的行块。质心力变换记 ${}^G\mathbf{X}_i^*$（对偶/力变换，对齐 \cref{ch:lie} 的伴随记号 $\mathrm{Ad}$）。
\end{note}

% ════════════════════════════════════════════════════════════════
\section{接触力的三大物理铁律}
\label{sec:contact-laws}

\paragraph{动机：无约束优化给出非物理解。}
设想一个四足机器人要向右加速，控制器解一个最小化关节力矩的二次规划，只把牛顿--欧拉方程 $\mathbf{M}\ddot{\boldsymbol{q}}+\boldsymbol{h}=\mathbf{S}^\top\boldsymbol{\tau}+\mathbf{J}_c^\top\boldsymbol{\lambda}$ 当等式约束、却忘了约束接触力本身。求解器会怎样作弊？它会令法向力 $\lambda^z\to0$（省力）、令水平摩擦力 $\lambda^x$ 任意大（全靠摩擦推动）——这在数学上「最优」，在物理上却意味着摩擦力远超法向力乘以摩擦系数，机器人当场打滑。更极端地，它可以让脚在离地 $10\,$cm 处仍受到地面「支撑力」，或让地面「吸住」机器人。这正是为什么\emph{每一个}足式控制器都必须显式编码接触力约束。本节给出这些约束的物理与数学形式——它们是全章的地基。

\paragraph{历史。}
接触力学的数学框架可追溯到四位先驱。Coulomb（1785）在《简单机械理论》中提出干摩擦定律 $f_t\le\mu f_n$\cite{coulomb1785theorie}；Signorini（1933 年发表简短札记提出，1959 年课程完整形式化）把单边接触抽象为一个带歧义边界条件的弹性力学问题，即 Signorini 问题\cite{signorini1959questioni}；Fichera（1963--64）证明了 Signorini 问题解的存在唯一性，这一工作被视为变分不等式理论的诞生\cite{fichera1964problemi}；Moreau（1988）以非光滑分析与凸分析为框架，给出有限自由度动力学中单边接触与干摩擦的现代形式\cite{moreau1988unilateral}。

\subsection{铁律一：法向非负与 Signorini 条件}
\label{subsec:signorini}

\paragraph{物理直觉。}
地面只能「推」机器人，不能「拉」——足底与地面之间没有粘性连接。设接触点 $\boldsymbol{p}_c$ 到地面的有符号距离为 $\phi(\boldsymbol{p}_c)$，接触力法向分量为 $\lambda_c^n$，则地面只推不拉要求 $\lambda_c^n\ge0$。但仅此不够：完整的 Signorini（Hertz--Signorini--Moreau）条件是三联立。

\begin{definition}[Signorini 单边接触条件]
\label{def:signorini}
接触点的有符号距离 $\phi$ 与法向接触力 $\lambda_c^n$ 满足
\begin{equation}
  \phi(\boldsymbol{p}_c)\ge0,\qquad \lambda_c^n\ge0,\qquad \phi(\boldsymbol{p}_c)\cdot\lambda_c^n=0.
  \label{eq:signorini}
\end{equation}
三式分别表达\emph{不穿透}（脚不能穿过地面）、\emph{法向力非负}（地面只推不拉）、\emph{互补性}（离地则无力、有力则在地面上）。
\end{definition}

\paragraph{几何解读：L 形可行域。}
在 $(\phi,\lambda_c^n)$ 平面上，合法的点只能落在两条射线上——正 $\phi$ 轴（空中、无力）和正 $\lambda_c^n$ 轴（接触、有力），构成一个「L 形」可行域（\cref{fig:contact-Lshape}）。这好比理想二极管：电压（对应 $\phi$）与电流（对应 $\lambda_c^n$）不能同时为正，要么导通、要么截止。不同之处在于二极管的伏安关系是标量的，而接触还叠加了切向摩擦，使问题从一维扩展到三维。

\begin{figure}[ht]
\centering
\begin{tikzpicture}[>=Stealth, scale=1.0]
  \draw[->] (-0.3,0) -- (4.2,0) node[right]{$\phi$ (gap)};
  \draw[->] (0,-0.3) -- (0,3.4) node[above]{$\lambda_c^n$ (法向力)};
  % feasible set: two rays forming an L
  \draw[ultra thick, main] (0,0) -- (4.0,0);
  \draw[ultra thick, main] (0,0) -- (0,3.2);
  \fill (0,0) circle (1.6pt);
  \node[main, anchor=south west] at (2.4,0) {接触有力\ $\phi=0,\ \lambda^n>0$};
  \node[main, anchor=west, rotate=90] at (0,1.4) {离地无力\ $\phi>0,\ \lambda^n=0$};
  \node[anchor=north east, font=\small] at (0,0) {O};
  % forbidden quadrant shading
  \fill[red!8] (0.05,0.05) rectangle (3.9,3.0);
  \node[red!60!black, font=\small] at (2.2,2.2) {禁区：$\phi\!>\!0$ 且 $\lambda^n\!>\!0$};
\end{tikzpicture}
\caption{Signorini 条件的「L 形」可行域：合法的 $(\phi,\lambda_c^n)$ 只落在两条坐标半轴上，二者不能同时为正——这正是互补性 $\phi\cdot\lambda_c^n=0$ 的几何含义。}
\label{fig:contact-Lshape}
\end{figure}

\paragraph{变分不等式 / 法锥观点。}
当地面是平面、可行集为 $C=\{\phi\ge0\}$ 时，\cref{eq:signorini} 可等价改写为一个法锥包含关系：

\begin{proposition}[Signorini 的法锥形式]
\label{prop:signorini-normalcone}
设凸集 $C=\{\phi\in\mathbb{R}:\phi\ge0\}$，其在点 $\phi$ 处的法锥
$N_C(\phi)=\{\,d:\,d^\top(\phi'-\phi)\le0,\ \forall\phi'\in C\,\}$。则 \cref{eq:signorini} 等价于 $-\lambda_c^n\in N_C(\phi)$，亦即 $\lambda_c^n\in -N_C(\phi)$。
\end{proposition}
\begin{proof}
$C=[0,+\infty)$ 的法锥分两种情形：在内部 $\phi>0$ 处 $N_C(\phi)=\{0\}$，故 $\lambda_c^n=0$；在边界 $\phi=0$ 处 $N_C(0)=(-\infty,0]$，故 $-\lambda_c^n\le0$ 即 $\lambda_c^n\ge0$。两情形合起来恰是「$\phi>0\Rightarrow\lambda^n=0$ 且 $\phi=0\Rightarrow\lambda^n\ge0$」，与 \cref{eq:signorini} 的三联立等价。$\square$
\end{proof}

\noindent 这个观点重要，因为它把一条「非光滑逻辑约束」翻译成了凸分析中标准的\emph{变分不等式}：接触力是约束可行集法锥中的元素。后续把接触塞进优化器、以及与 KKT 互补松弛（\cref{ch:nlopt}）建立联系，都以此为桥。

\paragraph{非平坦地形的法向。}
对于非平坦地形，法向 $\boldsymbol{n}$ 不再是 $[0,0,1]^\top$。设地形高度函数 $z=h(x,y)$，则单位法向
\begin{equation}
  \boldsymbol{n}=\frac{1}{\sqrt{1+h_x^2+h_y^2}}\,[-h_x,\,-h_y,\,1]^\top,\qquad h_x=\partial h/\partial x,\ h_y=\partial h/\partial y,
  \label{eq:terrain-normal}
\end{equation}
法向非负约束变为 $\lambda_c^n=\boldsymbol{\lambda}_c\cdot\boldsymbol{n}\ge0$。在倾斜地面上仍用 $\lambda_c^z\ge0$ 是常见错误——它允许了沿斜面的「推力」却没有正确约束真实法向力，会导致机器人在斜坡上滑动。约束必须在\emph{接触面局部坐标系}中表达。

\subsection{铁律二：库仑摩擦锥}
\label{subsec:friction-cone}

\paragraph{物理直觉。}
摩擦力不能无限大：切向力的模不超过法向力乘以摩擦系数。

\begin{definition}[库仑摩擦锥]
\label{def:friction-cone}
切向接触力满足
\begin{equation}
  \|\boldsymbol{\lambda}_c^t\|_2=\sqrt{(\lambda_c^x)^2+(\lambda_c^y)^2}\le\mu\,\lambda_c^n,
  \label{eq:friction-cone-contact}
\end{equation}
其中 $\mu$ 是静摩擦系数。合法的接触力向量在 $(\lambda^x,\lambda^y,\lambda^z)$ 空间中落在一个圆锥内：顶点在原点、轴线沿法向 $z$、\textbf{半顶角} $\alpha=\arctan\mu$。
\end{definition}

\noindent 当 $\mu=0.7$ 时 $\alpha\approx35^\circ$；当 $\mu=0.1$（冰面）时 $\alpha\approx5.7^\circ$——锥非常「尖」，几乎只能产生竖直力。典型摩擦系数：橡胶--干水泥 $0.6\sim0.8$，橡胶--湿水泥 $0.4\sim0.5$，橡胶--冰 $0.05\sim0.15$，金属--金属 $0.3\sim0.6$。

\paragraph{为什么是二阶锥？}\cref{eq:friction-cone-contact} 在数学上是一个\emph{二阶锥（SOC）}约束。SOC 是\textbf{凸}的，但不是线性的：

\begin{proposition}[摩擦锥的凸性]
\label{prop:cone-convex}
\cref{eq:friction-cone-contact} 定义的集合是凸锥。
\end{proposition}
\begin{proof}
设 $\boldsymbol{\lambda}_1,\boldsymbol{\lambda}_2$ 都在锥内，$\|\boldsymbol{\lambda}_i^t\|\le\mu\lambda_i^n$。对任意 $\alpha\in[0,1]$，由三角不等式与各自的锥约束，
\[
\|\alpha\boldsymbol{\lambda}_1^t+(1-\alpha)\boldsymbol{\lambda}_2^t\|\le\alpha\|\boldsymbol{\lambda}_1^t\|+(1-\alpha)\|\boldsymbol{\lambda}_2^t\|\le\mu\big(\alpha\lambda_1^n+(1-\alpha)\lambda_2^n\big),
\]
故凸组合仍在锥内。又因约束对正标量缩放不变，是锥。$\square$
\end{proof}

\noindent 凸性的后果是分工：若问题只有 SOC 约束加二次目标，得到的是\emph{二阶锥规划（SOCP）}，可用 Mosek、ECOS 等专门的锥求解器；但若想用通用 QP 求解器（OSQP、qpOASES），就必须把 SOC \emph{线性化}成多面体——这正是 \cref{sec:contact-linearize} 的内容。

\subsection{铁律三：互补性}
\label{subsec:complementarity}

一只脚要么在空中（无接触力）、要么在地面上（有接触力），不可能「在空中却受地面力」。把这条逻辑写成数学：
\begin{equation}
  0\le\phi(\boldsymbol{p}_c)\ \perp\ \lambda_c^n\ge0,
  \label{eq:complementarity}
\end{equation}
符号 $\perp$ 表示 $\phi\cdot\lambda_c^n=0$——两个非负量乘积为零，意味着至少一个为零。若 $\phi$ 是 $\boldsymbol{q}$ 的线性函数（或线性化后），整个接触问题就成为线性互补问题（LCP，\cref{sec:contact-complementarity-lcp}）。

\paragraph{速度级与加速度级的 Signorini。}
位置级互补 $\phi\cdot\lambda^n=0$ 在时间积分中不够用，还需约束速度。设接触点法向相对速度 $\dot\phi$、法向相对加速度 $\ddot\phi$：
\begin{align}
  \text{分离：}\quad &\phi=0,\ \dot\phi>0\ \Rightarrow\ \lambda^n=0,\\
  \text{保持：}\quad &\phi=0,\ \dot\phi=0,\ \ddot\phi\ge0,\ \lambda^n\ge0,\ \ddot\phi\cdot\lambda^n=0.
  \label{eq:signorini-accel}
\end{align}

\paragraph{为什么互补约束难处理？}
互补约束是一个\emph{非光滑}的等式约束：在 $(\phi,\lambda^n)=(0,0)$（接触刚建立/断裂的瞬间）不可微，就像绝对值函数 $|x|$ 在 $x=0$ 处的棱角，只是这里棱角发生在二维。其后果是：标准梯度法在此处约束 Jacobian 退化（违反约束规格 LICQ），目标梯度含不连续跳变，传统 QP 求解器无法直接处理。这正是 \cref{sec:contact-routes} 三大路线要解决的核心难题。

\begin{pitfall}[概念误区：互补性 = 碰撞检测？]
互补约束描述的是\emph{稳态}接触关系——「接触 $\leftrightarrow$ 力」的逻辑耦合。碰撞（impact）是一个\emph{瞬时事件}，需要额外的冲量模型（如恢复系数），其数值处理（碰撞冲量、恢复系数、LCP 的 Lemke/PATH 数值求解）留待\cref{ch:constrained-dynamics} 的约束动力学讨论；本章只处理稳态互补性（足式 MPC/WBC 通常假设软着陆、无碰撞）。
\end{pitfall}

\begin{pitfall}[思维陷阱：三大约束彼此独立、互补性多余？]
若没有互补性，优化器可以在脚离地 $10\,$cm 时仍产生接触力——这在数学上满足 $\lambda^z\ge0$ 与摩擦锥，物理上却是胡说。互补性是连接\emph{几何（距离）}与\emph{力学（力）}的桥梁，缺一不可。三条铁律并非三条独立规则的叠加，而是同一个物理原理——接触是\emph{单边约束}（只推不拉、只在几何接触时产力）——的三个投影。
\end{pitfall}

\subsection{统一框架与最大耗散原则}
\label{subsec:max-dissipation}

把三大约束合写为紧凑形式：
\begin{equation}
  \begin{aligned}
  &\text{(Signorini)} && 0\le\phi(\boldsymbol{p}_c)\ \perp\ \lambda_c^n\ge0,\\
  &\text{(Coulomb)} && \|\boldsymbol{\lambda}_c^t\|_2\le\mu\,\lambda_c^n,\\
  &\text{(最大耗散)} && \dot{\boldsymbol{p}}_c^t\neq\boldsymbol{0}\ \Rightarrow\ \boldsymbol{\lambda}_c^t=-\mu\lambda_c^n\,\frac{\dot{\boldsymbol{p}}_c^t}{\|\dot{\boldsymbol{p}}_c^t\|}.
  \end{aligned}
  \label{eq:contact-unified}
\end{equation}
第三式是\textbf{最大耗散原则}：若接触点在滑动，摩擦力方向与滑动方向相反、大小恰为 $\mu\lambda_c^n$（锥边界）。这是库仑摩擦的完整表述，可由「摩擦耗散功率在所有满足摩擦锥的力中最大」导出，故名。多数 MPC/WBC 假设\emph{无滑动} $\dot{\boldsymbol{p}}_c^t=0$，此时第三式自动满足，只需处理前两式——这把问题从一般非凸的摩擦互补简化为「法向互补 + 凸摩擦锥」，是后续线性化与凸 MPC 的前提。

\paragraph{两类工程附加约束。}
除三大铁律外，工程中还有两类实用约束：\emph{足端力幅度}——摆动腿足端力严格为零，支撑腿法向力有上限 $f_{z,\max}$（电机物理极限），由步态调度器的接触状态变量 $c_i\in\{0,1\}$ 实现（$c_i=0$ 时强制 $\boldsymbol{f}_i=\boldsymbol{0}$）；\emph{关节力矩限幅}——接触力经 $\boldsymbol{\tau}=\mathbf{J}_c^\top\boldsymbol{\lambda}$ 映射到关节后须满足 $|\tau_j|\le\tau_{j,\max}$。当构型 $\boldsymbol{q}$ 也是决策变量时该约束非凸；但实时 MPC 中 $\boldsymbol{q}$ 由当前状态或参考轨迹确定，$\mathbf{J}_c$ 退化为常数矩阵，力矩约束变为关于足端力的线性约束，可直接纳入 QP。

\begin{practice}
\begin{enumerate}
  \item 一个四足机器人站在 $30^\circ$ 斜坡上。写出每只脚的法向非负约束与摩擦锥约束（用接触面局部系）。若 $\mu=0.5$，机器人能否稳定站立？（提示：比较重力相对斜面法向的倾角 $30^\circ$ 与摩擦锥半顶角 $\arctan0.5\approx26.6^\circ$；\cref{sec:contact-cone-giwc} 将给出严格的 LP 判据。）
  \item 用 \cref{prop:signorini-normalcone} 把 Signorini 三联立写成 $-\lambda^n\in N_C(\phi)$。说明这一变分不等式观点为何便于与 \cref{ch:nlopt} 的 KKT 互补松弛建立联系。
  \item 仿真器（如 MuJoCo）常用软接触（弹簧--阻尼）。软接触是否严格满足 Signorini 条件？为何仿真器愿意违反它？（\cref{sec:contact-soft} 将系统对比。）
\end{enumerate}
\end{practice}

% ════════════════════════════════════════════════════════════════
\section{摩擦锥线性化：从 SOC 到 QP}
\label{sec:contact-linearize}

\paragraph{承上启下。}
\cref{sec:contact-laws} 确立了摩擦锥是凸但非线性的 SOC 约束。通用 QP 求解器吃不下 SOC，于是 \cref{ch:control} 的全身控制把它「线性化成金字塔」。本节补齐那一笔背后的全部细节：多面体内/外切的推导、误差界、以及「内切才保守」这一安全性辨析——它们决定了你的机器人在摩擦极限附近会不会滑倒。

\subsection{多面体内/外切近似}
\label{subsec:polyhedral}

\paragraph{思想。}
用 $k$ 条边的正多边形内切或外切于圆，把圆锥近似为多面体锥——正如图形学用多边形逼近曲面（边数越多越圆），区别在于这里的逼近误差直接影响机器人是否打滑。在 $\lambda^z=1$ 平面上截取摩擦锥得半径为 $\mu$ 的圆 $(\lambda^x)^2+(\lambda^y)^2\le\mu^2$。在圆周上等距取 $k$ 点 $(\mu\cos\theta_i,\mu\sin\theta_i)$，$\theta_i=2\pi i/k$，连成正 $k$ 边形，每条边给一条线性不等式：第 $i$ 条要求接触力在外法向上的投影不超过边到原点的距离，

\begin{equation}
  \boxed{\ \cos\theta_i\,\lambda^x+\sin\theta_i\,\lambda^y\le\mu_{\mathrm{eff}}\,\lambda^z,\qquad i=0,1,\dots,k-1\ }
  \label{eq:cone-linear-ineq}
\end{equation}

\noindent（这是全章唯一保留的公式级「实现」形态：堆叠 $n_c$ 个接触、各 $k$ 条，即得 WBC-QP 的摩擦锥不等式块。）有效系数 $\mu_{\mathrm{eff}}$ 决定内/外切：
\begin{equation}
  \mu_{\mathrm{eff}}=\begin{cases} \mu, & \text{外切近似（outer）},\\[2pt] \mu\cos(\pi/k), & \text{内切近似（inner）}. \end{cases}
  \label{eq:mu-eff}
\end{equation}
内切的来历：正 $k$ 边形内切于半径 $\mu$ 的圆时，其\emph{内切圆半径}为 $r_{\mathrm{inner}}=\mu\cos(\pi/k)$（边中点到圆心的距离）。$k=4$ 时 $r_{\mathrm{inner}}=\mu/\sqrt2\approx0.707\mu$，损失约 $29.3\%$ 的摩擦能力；$k=8$ 时 $\approx0.924\mu$，损失 $7.6\%$；$k=16$ 时 $\approx0.981\mu$，损失 $1.9\%$。

\begin{figure}[ht]
\centering
\begin{tikzpicture}[scale=1.6, >=Stealth]
  \def\mu{1.0}
  % the friction circle (cross-section at lambda^z=1)
  \draw[main, thick] (0,0) circle (\mu);
  \node[main, anchor=south west] at (0.62,0.62) {真实锥截面（圆，半径 $\mu$）};
  % outer square (k=4, axis-aligned)  -> shows box-cone over-relaxation
  \draw[red!70!black, thick] (-\mu,-\mu) rectangle (\mu,\mu);
  % inner polygon k=4 (rotated 45 to be inscribed cleanly) -- inscribed square
  \draw[second, thick, fill=second!10, fill opacity=0.4]
    (\mu,0) -- (0,\mu) -- (-\mu,0) -- (0,-\mu) -- cycle;
  \node[second, anchor=north] at (0,-1.18) {内切（$k=4$）：$r=\mu\cos45^\circ$，保守};
  \node[red!70!black, anchor=south] at (1.0,1.02) {外切（$k=4$，轴对齐）：box cone，过松};
  \draw[->] (0,0) -- (\mu,0) node[midway, below, font=\small]{$\mu$};
  \fill (0,0) circle (0.6pt);
\end{tikzpicture}
\caption{摩擦锥在 $\lambda^z=1$ 平面上的截面与 $k=4$ 多面体近似。外切方形锥（红，轴对齐）在对角方向允许 $\sqrt2\,\mu\lambda^z$ 的摩擦力，过度放松；内切正方形（蓝）始终落在真实圆内，保守。}
\label{fig:cone-cross-section}
\end{figure}

\subsection{近似误差的定量分析}
\label{subsec:cone-error}

外切多面体比圆锥大多少，可用面积比衡量。外切正 $k$ 边形与圆的面积比
\begin{equation}
  \frac{A_{\mathrm{outer}}}{A_{\mathrm{circle}}}=\frac{k\tan(\pi/k)}{\pi},\qquad
  \frac{A_{\mathrm{inner}}}{A_{\mathrm{circle}}}=\frac{k\sin(2\pi/k)}{2\pi}.
  \label{eq:area-ratio}
\end{equation}
代入典型 $k$（\cref{tab:cone-error}），$k=4$ 外切超出 $27.3\%$、内切损失 $36.3\%$；$k=8$ 分别为 $5.5\%$ 与 $10.0\%$；$k=16$ 降到 $1.3\%$ 与 $2.5\%$。径向误差按 $1-\cos(\pi/k)$（内切）尺度衰减。

\begin{table}[H]\centering\small
\caption{摩擦锥多面体近似的面积误差（$k$ 为边数）}
\label{tab:cone-error}
\begin{tabular}{ccccc}
\toprule
$k$ & 外切面积比 & 外切超出 & 内切面积比 & 内切损失 \\
\midrule
4 & 1.273 & $27.3\%$ & 0.637 & $36.3\%$ \\
6 & 1.103 & $10.3\%$ & 0.827 & $17.3\%$ \\
8 & 1.055 & $5.5\%$  & 0.900 & $10.0\%$ \\
12& 1.023 & $2.3\%$  & 0.955 & $4.5\%$ \\
16& 1.013 & $1.3\%$  & 0.975 & $2.5\%$ \\
\bottomrule
\end{tabular}
\end{table}

\begin{pitfall}[概念误区：外切更「安全」？恰恰相反]
直觉以为外切包含整个真实锥所以更安全——错。\textbf{内切才保守}。外切允许的力方向\emph{超出}真实摩擦锥，优化器可能找到一个「合法」解，实际执行时却滑倒；内切保守地限制力方向，宁可少用摩擦也不冒险。Cheetah 用外切，是因其 MPC 代价函数已倾向于产生接近竖直的力（加权使 $\lambda^z$ 优先），解很少触碰锥边界；但若控制器频繁在锥边界操作（如极端地形），应切到内切。
\end{pitfall}

\begin{pitfall}[思维陷阱：$k=4$ 轴对齐退化为盒形锥]
取 $\theta_i=0,\tfrac\pi2,\pi,\tfrac{3\pi}2$（与 $x$-$y$ 轴对齐），$k=4$ 多面体退化成\emph{方形锥}（box cone），约束变成 $|\lambda^x|\le\mu\lambda^z$ 且 $|\lambda^y|\le\mu\lambda^z$，对角方向上竟允许 $\sqrt2\mu\lambda^z$ 的摩擦力（\cref{fig:cone-cross-section} 红框）——过度放松。解决：旋转 $45^\circ$ 取 $\theta_i=\tfrac\pi4,\tfrac{3\pi}4,\dots$，或直接用 $k=8$。
\end{pitfall}

\begin{insight}[伪精度：线性化误差与 $\mu$ 测不准同量级]
追求 $k=16$ 才安全往往是伪精度。摩擦系数 $\mu$ 本身的测量误差通常在 $10\%\sim20\%$，与 $k=4$ 的近似误差（约 $27\%$）、$k=8$ 的约 $5.5\%$ 是同一量级。在 $\mu$ 测不准时，更好的策略是\emph{用 $k=4$ 或 $k=8$ 配合保守的 $\mu$ 估计}（真实 $\mu=0.7$ 时控制器用 $\mu=0.5$）。精度--性能权衡的关键在边数与 $\mu$ 不确定性的匹配，而非盲目增大 $k$——这也是为什么实时凸 MPC 普遍取 $k=4$、标准 WBC 取 $k=8$，无须求精确 SOC（精确 SOC 需 Mosek/ECOS 等 SOCP 求解器；通用 QP 求解器如 OSQP、ProxQP 处理的是\emph{线性化后}的多面体摩擦锥）。
\end{insight}

\begin{practice}
\begin{enumerate}
  \item 由 \cref{eq:area-ratio} 推导 $k\to\infty$ 时两个面积比都趋于 $1$，并证明外切总比内切误差小一个 $\cos$ 因子（同一 $k$ 下）。
  \item 验证 $k=4$ 轴对齐时对角方向的最大摩擦力为 $\sqrt2\mu\lambda^z$；旋转 $45^\circ$ 后这一过松还在吗？
\end{enumerate}
\end{practice}

% ════════════════════════════════════════════════════════════════
\section{锥的两种表示与 GIWC 静平衡判据}
\label{sec:contact-cone-giwc}

\paragraph{承上启下。}
\cref{sec:contact-linearize} 把摩擦锥写成「一堆线性不等式」$\mathbf{A}\boldsymbol{\lambda}\le0$。这只是多面体锥的一种表示。本节先把它推广到\emph{面接触}的 6D wrench cone，再揭示锥的另一种等价表示（生成元），最后用二者把「机器人能否站稳」化为一个几何包含问题——这是本章的理论高峰之一，它把摩擦锥、ZMP 与支撑多边形统一了起来。

\subsection{Wrench Cone：从 3D 点接触到 6D 面接触}
\label{subsec:wrench-cone}

3D 点接触假设接触发生在一个几何点。但足底是有面积的矩形/圆形区域，面接触除了产生力，还产生\emph{力矩}（绕法向的扭转摩擦、压力中心偏移）。对一个矩形面接触（尺寸 $2l_x\times2l_y$），接触 wrench $\boldsymbol{w}_c=[\boldsymbol{f}_c;\boldsymbol{n}_c]\in\mathbb{R}^6$（本书序，力在前）须满足一组线性不等式 $\mathbf{U}\boldsymbol{w}_c\le0$，其中 $\mathbf{U}$ 同时编码：法向力非负 $f_c^z\ge0$；摩擦锥 $\|\boldsymbol{f}_c^{xy}\|\le\mu f_c^z$（线性化后）；绕法向的扭转摩擦 $|n_c^z|\le\gamma f_c^z$（$\gamma$ 为扭转摩擦系数）；压力中心（CoP）约束 $|n_c^x|\le l_y f_c^z$ 且 $|n_c^y|\le l_x f_c^z$（CoP 不能跑出足底）。

\begin{note}[排序对照：Caron 2015 的 $\mathbf{U}$ 矩阵须翻转行块]
Caron 等（2015）给出的单矩形面 CWC 闭式 $\mathbf{U}$ 矩阵\cite{caron2015stability} 用\emph{力矩在前} $[\boldsymbol{n}_c;\boldsymbol{f}_c]$ 排序，前 3 列对应 $n^{x,y,z}$、后 3 列对应 $f^{x,y,z}$。本书 wrench 取力在前 $[\boldsymbol{f}_c;\boldsymbol{n}_c]$，故须把其每一行的「力矩 3 列」与「力 3 列」对调（等价于右乘一个 $6\times6$ 置换 $\mathbf{P}$）。翻转后本书序的 $\mathbf{U}$ 形如（按上述四组约束逐行）：法向行 $[-1,0,0,\ 0,0,0]$；摩擦锥四行 $[\pm1,0,-\mu,\ \mathbf{0}_3]$、$[0,\pm1,-\mu,\ \mathbf{0}_3]$；CoP 四行 $[0,0,-l_y,\ \pm1,0,0]$、$[0,0,-l_x,\ 0,\pm1,0]$；扭转两行 $[0,0,-\gamma,\ 0,0,\pm1]$。数值结构与 Caron 一致，仅列序按本书 $[\boldsymbol{f};\boldsymbol{n}]$ 重排。
\end{note}

\noindent 关键在于\textbf{力与力矩不独立}：如 Caron 等指出，固定合力后，力矩的可行范围（moment cone）形状随之改变，故 6D wrench cone \emph{不是} 3D 力锥与 3D 力矩锥的笛卡尔积，二者有耦合约束。对人形扁平足，6D wrench cone 比 3D 点接触更准——它能捕捉「脚尖翘起」（CoP 移到足底边缘）这类失稳模式。

\subsection{\texorpdfstring{$H$}{H}-表示与 \texorpdfstring{$V$}{V}-表示的对偶}
\label{subsec:hrep-vrep}

到此为止，摩擦锥/wrench cone 都写成「半空间的交」$\mathbf{A}\boldsymbol{\lambda}\le0$，称\textbf{半空间表示}（$H$-rep）。同一个多面体锥还有完全等价的另一种表示——\textbf{生成元表示}（$V$-rep）：锥是若干「射线」（生成元）的非负组合。把 $k$ 边线性化摩擦锥的 $k$ 条棱写成生成元（接触局部系，$z$ 为法向）：
\begin{equation}
  \boldsymbol{g}_i=[\mu\cos\theta_i,\ \mu\sin\theta_i,\ 1]^\top,\quad \theta_i=\tfrac{2\pi i}{k},\quad
  \boldsymbol{\lambda}=\sum_{i=0}^{k-1}\beta_i\boldsymbol{g}_i,\ \ \beta_i\ge0.
  \label{eq:vrep-generators}
\end{equation}

\begin{theorem}[Minkowski--Weyl]
\label{thm:minkowski-weyl}
每个多面体锥既可写成有限多个半空间的交（$H$-rep $\mathbf{A}\boldsymbol{\lambda}\le0$），也可写成有限多条射线的非负包（$V$-rep $\boldsymbol{\lambda}=\sum_i\beta_i\boldsymbol{g}_i,\,\beta_i\ge0$），两者可相互转换。
\end{theorem}

\begin{insight}[$H$-rep 与 $V$-rep：同一个锥的「内」「外」两种视角]
$H$-rep 从外部「切」——每个不等式砍掉一个半空间，适合\emph{检验}一个给定的力是否合法（逐行代入即可）。$V$-rep 从内部「撑」——每条生成元从原点向外射，适合\emph{生成/参数化}所有合法的力（给一组 $\beta_i\ge0$ 即得一个合法力）。优化建模时这个区别是实打实的：WBC 把接触力当决策变量、要施加约束，故用 $H$-rep（$\mathbf{A}\boldsymbol{\lambda}\le0$ 直接进 QP）；而可达性分析要枚举所有可能的合力，故用 $V$-rep。
\end{insight}

\paragraph{双重描述法与组合爆炸。}
$V$-rep $\leftrightarrow$ $H$-rep 的经典转换算法是\textbf{双重描述法}（Double Description, DD；Motzkin 等 1953\cite{motzkin1953double}），`pycddlib` 是常用实现。当多个接触协同（人形双脚加双手、四足四脚），把各接触的力锥变换到公共参考系（质心）并求 Minkowski 和，得到的总锥即\textbf{广义惯性 wrench 锥}（Gravito-Inertial Wrench Cone, GIWC）。计算它有两条路：单矩形面用 Caron 2015（ICRA）的闭式 $\mathbf{U}$ 矩阵\cite{caron2015stability}；多接触用 Caron 2015（RSS）的双重描述法\cite{caron2015leveraging}——先把各接触锥的生成元（$V$-rep）合成，再转回 $H$-rep。

\begin{pitfall}[陷阱：$V$-rep 与 $H$-rep 的组合爆炸方向相反]
别以为「边数 $k$ 越大，两种表示都只线性增长」。对多接触 Minkowski 和，$V$-rep 的生成元数是各接触之\emph{和}（$n_c\cdot k$，线性）；但转成 $H$-rep 后的不等式行数可能\emph{指数级}膨胀（最坏 $O(k^{n_c})$，McMullen 上界定理），DD 在高维退化锥上会很慢。正确做法：多接触时\emph{离线}预计算 GIWC 的 $H$-rep（接触几何固定时锥不变），在线只做矩阵乘法判定；或直接在 WBC 里用每接触的 $H$-rep 而不显式构造 GIWC。
\end{pitfall}

\subsection{静平衡可行性：「机器人能站稳吗？」}
\label{subsec:static-equilibrium}

\paragraph{静平衡条件。}
机器人静止（$\dot{\boldsymbol{v}}=\boldsymbol{0}$，所有 link 速度为零）时，牛顿--欧拉方程退化为「接触力的合力/合力矩 = 重力的合力/合力矩」。把所有接触力变换到质心 $G$ 求和，记总接触 wrench $\boldsymbol{w}_G^{\mathrm{contact}}=\sum_c{}^G\mathbf{X}_c^*\boldsymbol{\lambda}_c$，则（本书序，线动量率部分在前）
\begin{equation}
  \boldsymbol{w}_G^{\mathrm{contact}}=-\boldsymbol{w}_G^{\mathrm{gravity}}=[-m\boldsymbol{g};\,\boldsymbol{0}]_G,
  \label{eq:static-balance}
\end{equation}
其中角动量率部分为 $\boldsymbol{0}$（重力过质心、绕质心力矩为零）。

\begin{theorem}[静平衡的 GIWC 判据]
\label{thm:giwc-lp}
机器人能静止站立 $\iff$ 重力 wrench 的相反数落在 GIWC 内：
\begin{equation}
  -\boldsymbol{w}_G^{\mathrm{gravity}}\in\mathcal{W}_{\mathrm{GIWC}}=\Big\{\sum_c{}^G\mathbf{X}_c^*\boldsymbol{\lambda}_c:\ \boldsymbol{\lambda}_c\in\text{摩擦锥}_c\Big\},
  \label{eq:giwc}
\end{equation}
等价于一个小型线性可行性问题（LP）：\emph{找} $\beta_{c,i}\ge0$ \emph{使} $\sum_{c,i}\beta_{c,i}\,({}^G\mathbf{X}_c^*\boldsymbol{g}_{c,i})=-\boldsymbol{w}_G^{\mathrm{gravity}}$。LP 可行则能站稳；不可行则无论怎样分配接触力都无法平衡重力，必滑倒或倾覆。
\end{theorem}

\noindent 这正是 $V$-rep 大显身手处：GIWC 是各接触摩擦锥（用生成元描述）变换、Minkowski 求和后张成的锥，判定「重力 wrench 是否在锥内」即上述 LP。

\paragraph{回到斜坡习题。}
四足站在 $30^\circ$ 斜坡、$\mu=0.5$。单脚摩擦锥半顶角 $\alpha=\arctan0.5\approx26.6^\circ$，重力相对斜面法向的倾角恰为坡角 $30^\circ$。由 $30^\circ>26.6^\circ$，重力方向落在摩擦锥之外——LP 无解，机器人无法仅靠静摩擦站稳，会下滑。临界摩擦系数 $\mu_{\min}=\tan30^\circ\approx0.577$。这就解答了 \cref{sec:contact-laws} 习题：$\mu=0.5$ 站不稳，需 $\mu\ge0.577$。

\begin{figure}[ht]
\centering
\begin{tikzpicture}[>=Stealth, scale=1.0]
  % slope
  \fill[gray!15] (0,0) -- (4.5,0) -- (4.5,2.6) -- cycle;
  \draw[thick] (0,0) -- (4.5,2.6);
  \node[anchor=north] at (2.2,-0.05) {$30^\circ$ 斜坡};
  % contact point + local normal
  \coordinate (C) at (2.2,1.27);
  \fill (C) circle (1.6pt);
  % slope direction angle = atan(2.6/4.5)=~30deg; normal perpendicular
  \def\nx{-0.5} \def\ny{0.866}  % normal rotated 30 from vertical
  \draw[->, second, thick] (C) -- ($(C)+(\nx,\ny)$) node[anchor=south west]{$\boldsymbol{n}$ (斜面法向)};
  % friction cone about the normal (half-angle 26.6)
  \draw[main, thick] (C) -- ($(C)+(0.08,1.06)$);
  \draw[main, thick] (C) -- ($(C)+(-0.95,0.55)$);
  \draw[main, dashed] ($(C)+(0.08,1.06)$) arc (85:150:1.06);
  \node[main, font=\small, anchor=south] at ($(C)+(-0.55,0.95)$) {摩擦锥 $\alpha\!=\!26.6^\circ$};
  % gravity (straight down) -- falls OUTSIDE the cone
  \draw[->, red!70!black, thick] (C) -- ($(C)+(0,-1.2)$) node[anchor=north]{$-m\boldsymbol{g}$};
  \draw[->, red!70!black, thick] (C) -- ($(C)+(0,1.2)$) node[anchor=south]{重力反向 $\notin$ 锥};
\end{tikzpicture}
\caption{斜坡静平衡判据：重力反向相对斜面法向偏 $30^\circ$，超出半顶角 $26.6^\circ$（$\mu=0.5$）的摩擦锥，落在锥外，GIWC 判据的 LP 无解，机器人下滑。临界 $\mu_{\min}=\tan30^\circ\approx0.577$。}
\label{fig:slope-giwc}
\end{figure}

\begin{insight}[力学问题化为几何包含；ZMP 是 GIWC 的退化投影]
静平衡判据把「力学问题」彻底化为「几何包含问题」——能不能站稳，等价于\emph{重力这一个向量是否落在接触摩擦锥张成的那个锥里}。这就是为什么足式领域如此痴迷各种「区域」（支撑多边形、ZMP 区域、质心静平衡区域）：它们都是 GIWC 在不同子空间上的\emph{投影}。ZMP 必须落在支撑多边形内，本质上就是「重力 + 惯性 wrench 落在 GIWC 内」在平地、点接触假设下的退化特例——这条线索把本章的摩擦锥与简化模型（\cref{ch:reduced-models}）的 ZMP 判据统一了起来。
\end{insight}

\begin{pitfall}[思维陷阱：静平衡可行 $\neq$ 动态稳定]
GIWC 判据只保证存在一组\emph{静态}接触力平衡重力（必要条件），它不考虑：关节力矩能否真的产生这组接触力（受力矩限幅）；受扰动后能否恢复（动态稳定性，需 capturability/MPC 分析）。「能摆出一个平衡姿态」与「能稳定维持」是两件事——前者是静力学、后者是控制。
\end{pitfall}

\begin{practice}
\begin{enumerate}
  \item 对人形扁平足底（$l_x=0.1\,$m, $l_y=0.05\,$m, $\mu=0.6$），按本节排序对照构造本书序的 6D wrench cone 矩阵 $\mathbf{U}$。验证 CoP 刚好在足底边缘（$n_c^x=l_y f_c^z$）时对应行恰取等号。
  \item 写出 \cref{thm:giwc-lp} 的 LP，并说明为什么把它写成 $V$-rep（生成元的非负组合）比 $H$-rep 更自然。
\end{enumerate}
\end{practice}

% ════════════════════════════════════════════════════════════════
\section{互补约束的处理（一）：从刚体动力学到 LCP}
\label{sec:contact-complementarity-lcp}

\paragraph{承上启下。}
\cref{subsec:complementarity} 断言「接触问题是 LCP 的标准形式」，却没展示\emph{刚体动力学方程怎么一步步变成 LCP}。这一步推导是理解后续三大路线的前提——只有看清 LCP 的结构，才能明白「预定义模式」如何绕开它、「软化」在做什么近似。本节补上这条链，它也是本章与\cref{ch:constrained-dynamics} 约束动力学的理论交界：LCP 的标准形与 Delassus 算子的\emph{定义}留在本章（它是 GIWC/WBC 的理论入口），LCP 的\emph{数值求解}（Lemke/PATH/PGS）与碰撞冲量留待约束动力学章。

\subsection{半隐式时间步与 Delassus 算子}
\label{subsec:lcp-derivation}

考虑速度级的半隐式时间步。浮动基动力学 $\mathbf{M}(\boldsymbol{q})\dot{\boldsymbol{v}}+\boldsymbol{h}=\mathbf{S}^\top\boldsymbol{\tau}+\mathbf{J}_c^\top\boldsymbol{\lambda}$，$\mathbf{M}$ 对称正定，接触 Jacobian 的法向行记 $\mathbf{J}_n$。求一个时间步 $\Delta t$ 后的接触冲量 $\boldsymbol{\lambda}$（已吸收 $\Delta t$，量纲为冲量 $\mathrm{N\cdot s}$）。

\paragraph{步骤一：离散动力学，解出步末速度。}
前向欧拉离散 $\dot{\boldsymbol{v}}\approx(\boldsymbol{v}^+-\boldsymbol{v})/\Delta t$：
\[
\mathbf{M}(\boldsymbol{v}^+-\boldsymbol{v})=\Delta t\,(\mathbf{S}^\top\boldsymbol{\tau}-\boldsymbol{h})+\mathbf{J}_c^\top\boldsymbol{\lambda}.
\]
合并已知量为\emph{自由速度} $\boldsymbol{v}^{\mathrm{free}}=\boldsymbol{v}+\Delta t\,\mathbf{M}^{-1}(\mathbf{S}^\top\boldsymbol{\tau}-\boldsymbol{h})$——即\emph{假设无接触力}时下一步会达到的速度。于是 $\boldsymbol{v}^+=\boldsymbol{v}^{\mathrm{free}}+\mathbf{M}^{-1}\mathbf{J}_c^\top\boldsymbol{\lambda}$。

\paragraph{步骤二：法向速度级互补。}
接触点法向相对速度 $\dot\phi^+=\mathbf{J}_n\boldsymbol{v}^+$。速度级 Signorini（\cref{eq:signorini-accel}）要求接触点要么分离（$\dot\phi^+>0$，无力 $\lambda^n=0$）、要么保持（$\dot\phi^+=0$，可有力 $\lambda^n\ge0$），二者互补：$0\le\dot\phi^+\perp\lambda^n\ge0$。为不穿透，加一个 Baumgarte 风格位置修正偏置 $b=\tfrac{\gamma}{\Delta t}\phi$（把当前微小穿透 $\phi<0$ 在一步内「推」回零，$\gamma\in(0,1]$），约束写成 $\dot\phi^++b\ge0$。

\paragraph{步骤三：Schur 补消去 $\boldsymbol{v}^+$。}
把步骤一的 $\boldsymbol{v}^+$ 代入 $\dot\phi^+=\mathbf{J}_n\boldsymbol{v}^+$：
\begin{equation}
  \dot\phi^+=\mathbf{J}_n\boldsymbol{v}^{\mathrm{free}}+\underbrace{\mathbf{J}_n\mathbf{M}^{-1}\mathbf{J}_n^\top}_{=\,\mathbf{A}}\,\lambda^n.
  \label{eq:delassus-substitution}
\end{equation}

\begin{definition}[Delassus 算子与标准 LCP]
\label{def:delassus}
$\mathbf{A}=\mathbf{J}_n\mathbf{M}^{-1}\mathbf{J}_n^\top$ 称为 \textbf{Delassus 算子}（接触空间的「等效逆惯量」，由 $\mathbf{M}^{-1}$ 这个 Schur 补把全身惯量投影到接触法向）。令常数项 $\boldsymbol{w}=\mathbf{J}_n\boldsymbol{v}^{\mathrm{free}}+b$，代入互补条件得标准线性互补问题 $\mathrm{LCP}(\mathbf{A},\boldsymbol{w})$：
\begin{equation}
  \boxed{\ 0\le\lambda^n\ \perp\ \mathbf{A}\lambda^n+\boldsymbol{w}\ge0\ }
  \label{eq:lcp-standard}
\end{equation}
即求 $\lambda^n\ge0$ 使 $\mathbf{A}\lambda^n+\boldsymbol{w}\ge0$ 且 $(\lambda^n)^\top(\mathbf{A}\lambda^n+\boldsymbol{w})=0$。
\end{definition}

\begin{insight}[LCP 是把两个缠绕逻辑用 Schur 补解开的必然结果]
LCP 不是凭空规定的「数学游戏」，而是把「接触力修正速度」（动力学）与「速度决定有无力」（单边约束）这两个互相缠绕的逻辑，用 Schur 补解开后的必然结果。Delassus 算子 $\mathbf{A}=\mathbf{J}_n\mathbf{M}^{-1}\mathbf{J}_n^\top$ 是这条推导的核心——它把「一个接触点的力如何通过整副身体的惯量影响所有接触点的速度」压缩成一个稠密的 $n_c\times n_c$ 矩阵。正是这个稠密耦合，使多接触问题不能逐点独立求解。
\end{insight}

\subsection{半正定性、可解性与摩擦的代价}
\label{subsec:lcp-psd}

\begin{proposition}[Delassus 算子半正定]
\label{prop:delassus-psd}
$\mathbf{A}=\mathbf{J}_n\mathbf{M}^{-1}\mathbf{J}_n^\top\succeq0$。
\end{proposition}
\begin{proof}
因 $\mathbf{M}\succ0$ 故 $\mathbf{M}^{-1}\succ0$，对任意 $\boldsymbol{x}$ 有 $\boldsymbol{x}^\top\mathbf{A}\boldsymbol{x}=(\mathbf{J}_n^\top\boldsymbol{x})^\top\mathbf{M}^{-1}(\mathbf{J}_n^\top\boldsymbol{x})\ge0$。$\square$
\end{proof}

\begin{theorem}[半正定 LCP 的可解性]
\label{thm:lcp-solvable}
当 $\mathbf{A}$ 对称半正定\emph{且问题可行}时，$\mathrm{LCP}(\mathbf{A},\boldsymbol{w})$ 等价于一个下有界凸 QP
\begin{equation}
  \min_{\lambda^n\ge0}\ \tfrac12(\lambda^n)^\top\mathbf{A}\lambda^n+\boldsymbol{w}^\top\lambda^n,
  \label{eq:lcp-qp}
\end{equation}
必有解（Cottle--Pang--Stone\cite{cottle1992linear}）。接触问题因 $\boldsymbol{w}=\mathbf{J}_n\boldsymbol{v}^{\mathrm{free}}+b$ 的构造天然可行。当 $\mathbf{J}_n$ 行满秩时 $\mathbf{A}\succ0$，解唯一。
\end{theorem}

\noindent 这里「可行」二字不可省：纯半正定本身\emph{不}无条件保证可解，需附可行性（或 copositive-plus 加约束规格）；在接触场景里 $\boldsymbol{w}$ 的构造使问题天然可行，故工程结论成立。\cref{eq:lcp-qp} 也解释了 \cref{ch:control} 全身控制里「在 $\mathbf{J}_c$ 不满秩时给 $\mathbf{J}_c\mathbf{M}^{-1}\mathbf{J}_c^\top$ 对角加 $\epsilon\mathbf{I}$ 阻尼以保证可逆」的工程做法——它把半正定退化情形正则化为严格正定。

\paragraph{加入摩擦后：混合 LCP。}
上面只处理法向。一旦把线性化摩擦锥（\cref{sec:contact-linearize} 的 $k$ 边）写进来，切向冲量表示为锥生成元的非负组合，滑动方向引入额外互补对，整个问题变成一个\emph{更大的混合 LCP}（Stewart--Trinkle 1996 经典格式\cite{stewart1996implicit}），维度从 $n_c$ 膨胀到约 $n_c(k+2)$。

\begin{pitfall}[概念误区：LCP 总是凸优化？]
当 $\mathbf{A}\succeq0$（对称半正定，纯法向或关联摩擦模型）时，LCP 等价于凸 QP \cref{eq:lcp-qp}，可用 QP 求解器。但\emph{加入库仑摩擦后} $\mathbf{A}$ 通常\emph{非对称}（Stewart--Trinkle 格式），对应的不再是凸 QP，而是需要 Lemke 算法或 PATH 求解器的一般 LCP。这就是为什么「刚体接触 + 真实库仑摩擦」在数学上比「纯法向接触」难一个量级。这些非对称 LCP 的数值求解归\cref{ch:constrained-dynamics}。
\end{pitfall}

% ════════════════════════════════════════════════════════════════
\section{互补约束的处理（二）：三大工程路线}
\label{sec:contact-routes}

看清 LCP 的结构后，回到 \cref{subsec:complementarity} 的核心难题：互补约束非光滑，传统求解器无法直接处理。工程上有三条路线。

\subsection{路线一：预定义接触模式（mode-scheduled）}
\label{subsec:route-mode}

最直接的一条路是\emph{根本不让求解器去判断接触状态}——步态是预先设计的，控制工程师\emph{告诉}求解器每条腿在每个时刻是接触还是摆动。设时域 $[0,T]$ 被分成 $M$ 个模式段，每段有一个接触模式 $\sigma_m\in\{0,1\}^{n_c}$（位掩码，四足 $n_c=4$）。在第 $m$ 段对第 $i$ 个接触点：$\sigma_m^{(i)}=1$（在地面）时令 $\phi_i=0,\ \lambda_i^n\ge0,\ \dot{\boldsymbol{p}}_{c,i}^t=0$（无滑动）；$\sigma_m^{(i)}=0$（离地）时令 $\phi_i\ge0,\ \lambda_i^n=0$，并赋摆腿轨迹。\textbf{互补性自动满足}：每段内 $\phi_i$ 与 $\lambda_i^n$ 中恰有一个被固定为零，乘积自然为零。

\begin{table}[H]\centering\small
\caption{预定义模式路线的优缺点}
\label{tab:route-mode}
\begin{tabular}{ll}
\toprule
维度 & 评价 \\
\midrule
实时性 & 每段内是凸 QP/NLP，可实时求解 \\
简单性 & 实现简单，工业界主流（凸 MPC、双线程 MPC 栈） \\
适应性 & 接触模式不是优化变量，不能自适应地形 \\
新颖步态 & 不能发现新步态（如用手支撑） \\
鲁棒性 & 依赖步态调度器的正确性 \\
\bottomrule
\end{tabular}
\end{table}

\noindent 这条路线正是 \cref{ch:control} Cheetah 凸 MPC 与下游腿足栈采用的范式；它与变分/PMP 章（\cref{ch:opt-variational-pmp}）的混合系统 PMP（stance/swing 切换）是同一「模式切换」思想的优化与控制两面。

\subsection{路线二：软化互补约束}
\label{subsec:route-relax}

用光滑函数\emph{近似}非光滑的互补条件，使问题变成标准 NLP。有三种典型做法。\textbf{松弛}：$\phi\cdot\lambda^n\le\epsilon$，允许二者同时为小正数，优化中逐步 $\epsilon\to0$。\textbf{对数障碍}：$\min\ \mathrm{cost}-\kappa(\ln\phi+\ln\lambda^n)+\tfrac{1}{2\rho}(\phi\lambda^n)^2$，障碍项保证 $\phi,\lambda^n>0$、惩罚项推乘积趋零，逐步减小 $\kappa,\rho$。\textbf{Fischer--Burmeister（FB）函数}——最优雅的一种：

\begin{definition}[Fischer--Burmeister NCP 函数]
\label{def:fb}
$\Phi_{\mathrm{FB}}(a,b)=\sqrt{a^2+b^2}-a-b$。
\end{definition}

\begin{theorem}[FB 零点等价于互补]
\label{thm:fb-zero}
$\Phi_{\mathrm{FB}}(a,b)=0\iff a\ge0,\ b\ge0,\ ab=0$。
\end{theorem}
\begin{proof}
逐情形：若 $a>0,b=0$，$\sqrt{a^2}-a=|a|-a=0$；若 $a=0,b>0$ 同理为 $0$。若 $a>0,b>0$，由 $(a+b)^2=a^2+2ab+b^2>a^2+b^2$ 得 $\sqrt{a^2+b^2}<a+b$，故 $\Phi_{\mathrm{FB}}<0\neq0$。若 $a<0$（或 $b<0$），$\sqrt{a^2+b^2}\ge|a|>a$，故 $\Phi_{\mathrm{FB}}>0\neq0$。综上 $\Phi_{\mathrm{FB}}=0$ 精确等价于互补条件。$\square$
\end{proof}

\noindent 于是把互补约束替换为等式 $\Phi_{\mathrm{FB}}(\phi,\lambda^n)=0$，它在 $a>0$ 或 $b>0$ 时可微，不可微点只在 $(0,0)$，可用\emph{正则化}处理：$\Phi_{\mathrm{FB}}^{\epsilon}(a,b)=\sqrt{a^2+b^2+\epsilon^2}-a-b$ 处处可微，$\epsilon\to0$ 时逼近原 FB。判定一对 $(a,b)$ 是否（近似）互补，只需算一个标量：

\begin{codebox}[Python]{正则化 Fischer--Burmeister 互补判定（示意）}
phi_fb = sqrt(a*a + b*b + eps*eps) - a - b   # regularized FB value
return abs(phi_fb) < tol                      # True: 0 <= a _|_ b >= 0
\end{codebox}

\begin{table}[H]\centering\small
\caption{三种软化方法对比}
\label{tab:softening}
\begin{tabular}{lcccc}
\toprule
方法 & 光滑性 & 收敛速度 & 精度 & 实现复杂度 \\
\midrule
松弛 & 光滑 & 慢（线性）& $O(\epsilon)$ & 低 \\
对数障碍 & 光滑 & 中等 & $O(\kappa)$ & 中 \\
Fischer--Burmeister & 半光滑 & 快（超线性）& 精确 & 高 \\
\bottomrule
\end{tabular}
\end{table}

\begin{pitfall}[陷阱：FB 函数「是光滑的」可直接喂 Ipopt？]
原始 FB 在 $(0,0)$ 处不可微（半光滑）。Ipopt 需要二阶导数，而 FB 在 $(0,0)$ 附近的 Hessian 不存在——须用正则化版 $\Phi_{\mathrm{FB}}^{\epsilon}$ 或专门的半光滑牛顿法。另一个常见错误是一开始就设 $\epsilon=10^{-8}$：问题条件数极差，优化器不收敛、NaN 爆炸。正确做法是从 $\epsilon=0.1$ 起逐步减半（典型调度 $\{0.1,0.05,0.01,0.005,0.001\}$）。
\end{pitfall}

\subsection{路线三：接触隐式轨迹优化（CITO）}
\label{subsec:route-cito}

把接触位置、时刻和力\emph{全部}作为决策变量，让优化器自主发现最优接触策略。标准 CITO 公式（Posa 等 2014\cite{posa2014direct}）是一个带互补约束的数学规划（MPCC）：
\begin{equation}
  \begin{aligned}
  \min_{\boldsymbol{x},\boldsymbol{u},\boldsymbol{\lambda}}\ & \sum_{k=0}^{N-1} l(\boldsymbol{x}_k,\boldsymbol{u}_k)+l_f(\boldsymbol{x}_N)\\
  \text{s.t.}\ & \boldsymbol{x}_{k+1}=f(\boldsymbol{x}_k,\boldsymbol{u}_k,\boldsymbol{\lambda}_k), && \text{(动力学)}\\
  & 0\le\phi(\boldsymbol{x}_k)\ \perp\ \lambda_k^n\ge0, && \text{(法向互补)}\\
  & \|\boldsymbol{\lambda}_k^t\|\le\mu\lambda_k^n, && \text{(摩擦锥)}\\
  & 0\le(\mu\lambda_k^n-\|\boldsymbol{\lambda}_k^t\|)\ \perp\ \|\dot{\boldsymbol{p}}_c^t\|\ge0. && \text{(滑动互补)}
  \end{aligned}
  \label{eq:cito}
\end{equation}
最后一条是\emph{滑动互补}：若摩擦力未饱和（在锥内部），则接触点不滑动。

CITO 的根本困难有四：\emph{组合爆炸}（$n_c$ 接触、$N$ 步，最多 $2^{n_c N}$ 种接触模式序列）；\emph{非凸非光滑}（互补约束使可行域非凸）；\emph{局部极小}（不同接触策略对应不同局部最优）；\emph{计算代价}（典型 CITO 需分钟到小时级求解）。注意 MPCC 在互补点正违反 LICQ（\cref{ch:nlopt}），这正是 \cref{subsec:complementarity} 所说的非光滑难题在轨迹优化层的体现。

\begin{table}[H]\centering\small
\caption{接触隐式轨迹优化代表性进展（截至成稿）}
\label{tab:cito-progress}
\begin{tabular}{lll}
\toprule
工作 & 核心思想 & 性质 \\
\midrule
Posa 等 2014\cite{posa2014direct} & MPCC 标准形（直接法）& CITO 奠基 \\
Manchester--Kuindersma 2020\cite{manchester2020variational} & 变分接触隐式 TO & 经典变分形式 \\
Pang 等 2023\cite{pang2023global} & 局部光滑化 + quasi-dynamic & 操作任务全局规划 \\
Kurtz 等（预印本 2023, IJRR 2025）\cite{kurtz2025inverse} & 逆动力学 CITO-MPC（IDTO）& 见正文说明 \\
\bottomrule
\end{tabular}
\end{table}

\begin{pitfall}[思维陷阱：CITO 能取代步态设计、能实时跑腿足吗？]
目前 CITO 在腿足上仍是\emph{离线工具}——用于发现新步态或验证概念，实时控制仍依赖预定义步态加 MPC/WBC，因为 CITO 的求解时间（秒级）远超实时要求（毫秒级）。需澄清一个常见误传：Kurtz 等的 IDTO 框架（预印本 2023、IJRR 2025）在 Drake 中实现实时 CITO-MPC，硬件上演示 $>100\,$Hz，但\emph{演示对象是 20 自由度双臂操作，并非腿足}；腿足实时 CITO 仍属开放问题。更晚近的工作（如 GPU 加速的连续时间逐次凸化）据其报告在 GPU 上较既有逐次凸化方法有数量级加速，但「最快/SOTA」属作者自报、未经独立复现共识，本书不强引具体数字。\rebuilt{CITO 实时化前沿（2025--2026 新作）联网仍存疑，待核。}
\end{pitfall}

% ════════════════════════════════════════════════════════════════
\section{接触 Jacobian：从速度到广义力}
\label{sec:contact-jacobian}

\paragraph{承上启下。}
前三节解决了「接触力应满足什么约束」与「如何在优化中处理这些约束」。但还有一个环节未打通：接触力作用在足端，而动力学方程写在关节空间——两者的桥梁是接触 Jacobian。

\subsection{对偶关系：虚功原理}
\label{subsec:jacobian-duality}

接触点 $\boldsymbol{p}_c=\boldsymbol{p}_c(\boldsymbol{q})$ 是广义坐标的函数，链式求导得 $\dot{\boldsymbol{p}}_c=\tfrac{\partial\boldsymbol{p}_c}{\partial\boldsymbol{q}}\dot{\boldsymbol{q}}=\mathbf{J}_c(\boldsymbol{q})\dot{\boldsymbol{q}}$，其中 $\mathbf{J}_c\in\mathbb{R}^{3\times n_v}$ 是接触 Jacobian 的线速度部分（关心角速度则用 $6\times n_v$ 的 6D Jacobian）。

\begin{theorem}[接触力的对偶映射]
\label{thm:jacobian-dual}
由虚功原理，接触力 $\boldsymbol{\lambda}_c$ 在虚位移 $\delta\boldsymbol{p}_c$ 上的虚功等于对应广义力在虚广义位移上的虚功：
$\boldsymbol{\lambda}_c^\top\delta\boldsymbol{p}_c=\boldsymbol{\lambda}_c^\top\mathbf{J}_c\delta\boldsymbol{q}=(\mathbf{J}_c^\top\boldsymbol{\lambda}_c)^\top\delta\boldsymbol{q}$，故接触力在关节空间等价于
\begin{equation}
  \boldsymbol{\tau}_{\mathrm{contact}}=\mathbf{J}_c^\top\boldsymbol{\lambda}_c.
  \label{eq:jacobian-transpose}
\end{equation}
\end{theorem}

\noindent 这正是动力学方程 $\mathbf{M}\ddot{\boldsymbol{q}}+\boldsymbol{h}=\mathbf{S}^\top\boldsymbol{\tau}+\mathbf{J}_c^\top\boldsymbol{\lambda}_c$ 中接触力项的来源。

\begin{insight}[运动与力共用同一映射的转置]
$\mathbf{J}_c$ 与 $\mathbf{J}_c^\top$ 构成一对\emph{对偶映射}——$\mathbf{J}_c$ 把关节速度「翻译」成接触点速度（运动方向），$\mathbf{J}_c^\top$ 把接触力「翻译」成关节广义力（力方向）。这不是巧合，而是虚功原理的必然结果：能量守恒要求运动与力的传递使用\emph{同一个线性映射的转置}。这个对偶关系贯穿整个机器人学——从末端执行器力控到 WBC 的约束力分配（\cref{ch:operational-space,ch:force-wbc}），都依赖于此。
\end{insight}

\subsection{漂移项与 Baumgarte 稳定化}
\label{subsec:jacobian-drift}

在约束接触（脚固定在地面）情形下，加速度级约束为 $\ddot{\boldsymbol{p}}_c=\mathbf{J}_c\ddot{\boldsymbol{q}}+\dot{\mathbf{J}}_c\dot{\boldsymbol{q}}=0$，即 $\mathbf{J}_c\ddot{\boldsymbol{q}}=-\dot{\mathbf{J}}_c\dot{\boldsymbol{q}}$。\cref{ch:control} 的全身控制已直接写下这条约束 $\mathbf{J}_c\dot{\boldsymbol{\nu}}=-\dot{\mathbf{J}}_c\boldsymbol{\nu}$，本节补的是\emph{为何需要漂移项及如何稳定化}。其中 $\dot{\mathbf{J}}_c\dot{\boldsymbol{q}}$ 称\textbf{漂移项}（bias acceleration）：若忽略它，约束加速度 $\ddot{\boldsymbol{p}}_c\neq0$，接触点会逐渐漂离地面（仿真中表现为脚缓慢滑动、数秒后偏移数厘米）。

即便正确计算了漂移项，数值积分仍可能累积微小约束漂移。\textbf{Baumgarte 稳定化}在加速度约束中加入位置与速度修正：
\begin{equation}
  \mathbf{J}_c\ddot{\boldsymbol{q}}+\dot{\mathbf{J}}_c\dot{\boldsymbol{q}}+2\alpha\dot{\boldsymbol{p}}_c+\beta^2(\boldsymbol{p}_c-\boldsymbol{p}_c^{\mathrm{ref}})=0,
  \label{eq:baumgarte}
\end{equation}
其中 $\alpha,\beta>0$ 为稳定化增益（典型 $\alpha=\beta=50$），把约束违反当作一个临界阻尼二阶系统拉回零。注意此处 $\beta$ 是 Baumgarte 增益，与 \cref{subsec:lcp-derivation} 的位置修正比例 $\gamma$ 是不同记号、不可混。

\subsection{参考系选择：为何 LWA 适合接触}
\label{subsec:jacobian-frame}

接触 Jacobian 可在不同参考系下表达，常见三种：\texttt{LOCAL}（速度在 frame 自身坐标系，Jacobian 稀疏但摩擦锥方向随 frame 旋转）；\texttt{WORLD}（速度在世界系、参考点在世界原点，含远距离叉乘项 $\boldsymbol{\omega}\times\boldsymbol{r}$，数值条件差）；\texttt{LOCAL\_WORLD\_ALIGNED}（LWA，速度在世界系朝向下表达、参考点在 frame 处）。三者关系 $\mathbf{J}_{\mathrm{LWA}}=\mathrm{diag}(\mathbf{R}_{wf},\mathbf{R}_{wf})\,\mathbf{J}_{\mathrm{LOCAL}}$，其中 $\mathbf{R}_{wf}\in\SO(3)$ 是 frame 在世界系的姿态（\cref{ch:lie}；面接触法向随姿态变化时须用右扰动约定）。

\textbf{为何 LWA 最适合接触力学}：摩擦锥 $\|\boldsymbol{\lambda}^{xy}\|\le\mu\lambda^z$ 假设 $z$ 轴是接触面法向。在世界系中（水平地面）法向就是 $z$ 轴，直接匹配；若用 \texttt{LOCAL}，法向随 frame 旋转，每次都要旋转摩擦锥约束矩阵，增加复杂性且易错。LWA 兼顾了 \texttt{LOCAL} 的稀疏性（参考点在 frame 处，避免远距离叉乘）与世界系的直觉性（摩擦锥自然表达）。

\begin{practice}
\begin{enumerate}
  \item 对一个固定基座的 3 自由度平面机械臂，手动推导接触 Jacobian $\mathbf{J}_c(\boldsymbol{q})$，验证 $\mathbf{J}_c^\top\boldsymbol{\lambda}$ 给出的关节力矩与直接力矩平衡的结果一致。
  \item 解释为何忽略漂移项会导致仿真中脚滑动；写出加 Baumgarte 后约束违反满足的二阶常微分方程，并说明 $\alpha=\beta$ 对应临界阻尼。
\end{enumerate}
\end{practice}

% ════════════════════════════════════════════════════════════════
\section{质心动量矩阵 CMM 与 CCRBA}
\label{sec:contact-cmm}

\paragraph{承上启下。}
\cref{sec:contact-jacobian} 把接触力接回了关节空间。另一座桥是把全身运动接到\emph{质心}动量——质心动力学 $\dot{\boldsymbol{h}}_G=\sum(\text{外部 wrench})$ 是简化模型 MPC（\cref{ch:reduced-models}）的骨架，而其中的 $\boldsymbol{h}_G$ 如何从关节状态高效算出，正是本节的质心动量矩阵（CMM）与 CCRBA 算法要解决的。朴素逐 link 累加是 $O(N^2)$，CCRBA 把它降到 $O(N)$。

\subsection{CMM 的定义与分块结构}
\label{subsec:cmm-def}

\begin{definition}[质心动量与 CMM]
\label{def:cmm}
系统在质心 $G$ 处的 6D 空间动量是各 link 空间动量之和
\begin{equation}
  \boldsymbol{h}_G=[\boldsymbol{l}_G;\,\boldsymbol{k}_G]=\sum_{i=1}^{N_b}{}^G\mathbf{X}_i^*\,\mathbf{I}_i\,\boldsymbol{v}_i,
  \label{eq:centroidal-momentum}
\end{equation}
其中 $\boldsymbol{l}_G=m\dot{\boldsymbol{p}}_G$ 是线动量（本书序在前）、$\boldsymbol{k}_G$ 是绕质心角动量，${}^G\mathbf{X}_i^*$ 是 link $i$ 到质心的空间力变换、$\mathbf{I}_i$ 是 link 空间惯量、$\boldsymbol{v}_i$ 是 link 空间速度。由于 $\boldsymbol{v}_i$ 是 $\dot{\boldsymbol{q}}$ 的线性函数，$\boldsymbol{h}_G$ 也是：
\begin{equation}
  \boldsymbol{h}_G=\mathbf{A}_G(\boldsymbol{q})\,\dot{\boldsymbol{q}},\qquad \mathbf{A}_G\in\mathbb{R}^{6\times n_v},
  \label{eq:cmm-linear}
\end{equation}
$\mathbf{A}_G$ 即\textbf{质心动量矩阵}（CMM），依赖构型 $\boldsymbol{q}$ 但不依赖速度。
\end{definition}

\noindent CMM 有三条关键性质。其一，\emph{线动量子块恒为 $m\mathbf{I}_3$}：因 $\boldsymbol{l}_G=m\dot{\boldsymbol{p}}_G$ 且质心速度对基座线速度的偏导是 $m\mathbf{I}_3$，与构型无关。其二，\textbf{质心复合刚体惯量}（CCRBI）
\begin{equation}
  \mathbf{I}_G=\mathbf{A}_G\mathbf{M}^{-1}\mathbf{A}_G^\top\in\mathbb{R}^{6\times6},
  \label{eq:ccrbi}
\end{equation}
描述「整个机器人被焊成一个刚体时在质心处的空间惯量」，随构型 $\boldsymbol{q}$ 变化。其三，CMM 与全身质量矩阵的关系 $\mathbf{A}_G(\boldsymbol{q})=\sum_{i=1}^{N_b}{}^G\mathbf{X}_i^*\,\mathbf{I}_i\,\mathbf{J}_i(\boldsymbol{q})$，其中 $\mathbf{J}_i$ 是 link $i$ 质心的空间 Jacobian——注意 $\mathbf{A}_G$ 是 $6\times n_v$（关节空间到 6D 动量空间），而 $\mathbf{M}$ 是 $n_v\times n_v$。

\subsection{CCRBA 算法}
\label{subsec:ccrba}

CCRBA（质心复合刚体算法）是 Orin、Goswami、Lee（2013）提出的 $O(N)$ 算法\cite{orin2013centroidal}。核心思想：\emph{从叶到根递推复合惯量，再从根到叶计算 CMM 的列}。这好比算一棵组织树的总薪资——先从基层向上逐级汇总（叶到根）得每部门人力成本，再从总部向下分配（根到叶）；CCRBA 用同样逻辑汇总空间惯量，避免逐 link 独立计算的 $O(N^2)$ 重复遍历。算法分五阶段：

\begin{algo}[CCRBA（Orin 等 2013）骨架]
\label{algo:ccrba}
\emph{相位 1（根到叶，正向运动学）}：对 $i=1\dots N_b$，由关节变换算 link $i$ 的 frame 变换 $\mathbf{X}_i$ 与空间速度 $\boldsymbol{v}_i=\mathbf{X}_i\boldsymbol{v}_{\mathrm{parent}(i)}+\mathbf{S}_i\dot{q}_i$。
\emph{相位 2（叶到根，复合惯量，同 CRBA）}：对 $i=N_b\downarrow1$，$\mathbf{I}_c(i)=\mathbf{I}(i)+\sum_{j\in\mathrm{child}(i)}\mathbf{X}_j^*\mathbf{I}_c(j)\mathbf{X}_j^{-1}$。
\emph{相位 3（质心与质心变换）}：$\boldsymbol{p}_G=\tfrac1m\sum_i m_i\boldsymbol{p}_i$，构造 ${}^G\mathbf{X}$。
\emph{相位 4（CMM 的列，根到叶）}：$\mathbf{I}_G={}^G\mathbf{X}^*\mathbf{I}_c(\mathrm{root})\,{}^G\mathbf{X}^{-1}$；对 $i=1\dots N_b$，$\mathbf{A}_G(:,i)={}^G\mathbf{X}^*\mathbf{I}_c(i)\mathbf{S}_i$。
\emph{相位 5（质心动量）}：$\boldsymbol{h}_G=\mathbf{A}_G\boldsymbol{v}$。
\end{algo}

\noindent 相位 1--4 各为 $O(N)$，相位 5 是 $O(N\cdot n_v)$ 的矩阵-向量乘；整体以 $O(N)$ 递推构造 $\mathbf{A}_G$。对 $N=30$（$n_v\sim36$）的人形，相对朴素 $O(N^2)$ 法加速约 $30$ 倍。

\begin{pitfall}[概念误区：CMM 就是质量矩阵 $\mathbf{M}$ 的前 6 行？]
$\mathbf{A}_G$ 把 $\dot{\boldsymbol{q}}$ 映射到质心处的 6D 动量；$\mathbf{M}$ 把 $\ddot{\boldsymbol{q}}$ 映射到关节空间广义力。两者维度不同（$\mathbf{A}_G$ 是 $6\times n_v$，$\mathbf{M}$ 是 $n_v\times n_v$）。$\mathbf{A}_G$ 涉及到质心的坐标变换，\emph{不是}简单截取 $\mathbf{M}$ 的某些行。CMM 与 $\mathbf{M}$ 的真实关系经 ${}^G\mathbf{X}^*\mathbf{I}_c$ 这一串质心力变换给出（见 \cref{algo:ccrba} 相位 4），属 Orin--Goswami--Lee 2013 的核心结果。
\end{pitfall}

\begin{pitfall}[思维陷阱：CCRBI 是常数？]
系统的空间惯量\emph{随构型变化}。四足蜷缩四肢时绕质心转动惯量远小于伸展时——这就是 CCRBI $\mathbf{I}_G$ 依赖 $\boldsymbol{q}$ 的原因，与花样滑冰运动员收臂加速旋转同理。这一点在 \cref{sec:contact-angular} 的角动量可控性中至关重要。
\end{pitfall}

\begin{practice}
\begin{enumerate}
  \item 验证 CMM 的线动量子块为何\emph{恒}等于 $m\mathbf{I}_3$（与 $\boldsymbol{q}$ 无关），从 $\boldsymbol{l}_G=m\dot{\boldsymbol{p}}_G$ 出发。
  \item 论证四条腿完全伸直与完全蜷缩两种构型下 CCRBI 角动量子块（左上 $3\times3$）的差异方向，给出物理解释。
\end{enumerate}
\end{practice}

% ════════════════════════════════════════════════════════════════
\section{角动量的特殊地位：守恒与可控性}
\label{sec:contact-angular}

\paragraph{承上启下。}
\cref{sec:contact-cmm} 给出了高效计算 6D 质心动量的工具。其中线动量 $\boldsymbol{l}_G=m\dot{\boldsymbol{p}}_G$ 的行为直观，但角动量 $\boldsymbol{k}_G$ 蕴含更微妙、对控制至关重要的性质，值得专门展开。

\subsection{角动量动力学}
\label{subsec:angular-dynamics}

线动量与角动量的演化方程：
\begin{equation}
  \dot{\boldsymbol{l}}_G=m\boldsymbol{g}+\sum_c\boldsymbol{f}_c,\qquad
  \dot{\boldsymbol{k}}_G=\sum_c(\boldsymbol{p}_c-\boldsymbol{p}_G)\times\boldsymbol{f}_c.
  \label{eq:momentum-rate}
\end{equation}
关键区别：所有外力直接叠加改变线动量，但只有相对质心有「力臂」的力才产生角动量力矩；\textbf{重力对角动量的力矩为零}（重力过质心，力臂为零）。无接触时 $\dot{\boldsymbol{l}}_G=m\boldsymbol{g}$（自由落体），而 $\dot{\boldsymbol{k}}_G=\boldsymbol{0}$（守恒）；单点接触时线动量可自由改变（改力大小），角动量力矩方向受限（只能绕接触点）。

\subsection{飞行相位的角动量守恒}
\label{subsec:flight-conservation}

\begin{theorem}[飞行相位角动量守恒]
\label{thm:angular-conservation}
机器人完全离地（飞行相位）时，重力绕质心力矩为零，故 $\dot{\boldsymbol{k}}_G=\boldsymbol{0}$，$\boldsymbol{k}_G=\text{const}$。
\end{theorem}

\noindent 两个推论。其一，\emph{跳跃中不能空中转身}：若起跳时偏航角动量 $k_G^z=0$，整个飞行阶段 $k_G^z=0$，机器人不可能在空中改变偏航——若落地需朝向不同方向，必须在\emph{起跳前}注入足够角动量。其二，\emph{后空翻需精确起跳角动量}：要在飞行时间 $T$ 内恰好绕俯仰轴旋转 $2\pi$，需 $k_G^y=I_{yy}\cdot\tfrac{2\pi}{T}$，其中 $T=2v_z/g$ 由起跳竖直速度决定，而 $I_{yy}$ 飞行中可通过蜷缩/伸展改变——这就是体操运动员的技巧。

\subsection{猫翻身悖论}
\label{subsec:cat-paradox}

\paragraph{悖论。}猫从倒置自由落下，能在空中翻转 $180^\circ$ 平稳落地。但飞行中角动量守恒、初始 $\boldsymbol{k}_G=0$，怎么可能旋转？

\paragraph{解答。}猫\emph{不是作为刚体旋转}，而是通过关节内部运动改变姿态：前半身蜷缩（减小转动惯量）快速旋转 $\Delta\theta_{\mathrm{front}}$，后半身伸展（增大转动惯量）反向慢速旋转 $\Delta\theta_{\mathrm{rear}}$。零总角动量约束
\begin{equation}
  k_G=I_{\mathrm{front}}\dot\theta_{\mathrm{front}}+I_{\mathrm{rear}}\dot\theta_{\mathrm{rear}}=0\ \Rightarrow\ \dot\theta_{\mathrm{rear}}=-\frac{I_{\mathrm{front}}}{I_{\mathrm{rear}}}\dot\theta_{\mathrm{front}}.
  \label{eq:cat}
\end{equation}
由 $I_{\mathrm{front}}\neq I_{\mathrm{rear}}$（不对称蜷缩），$|\dot\theta_{\mathrm{front}}|\neq|\dot\theta_{\mathrm{rear}}|$，每一拍前后半身的转角不相互抵消，\emph{净旋转角非零}；重复多次，净旋转逐步累积。

\begin{figure}[ht]
\centering
\begin{tikzpicture}[>=Stealth, scale=1.0]
  % front body: curled (small inertia), fast rotation
  \draw[main, thick, fill=main!10] (0,0) circle (0.45);
  \node[main, font=\small] at (0,0) {前};
  \draw[->, main, thick] (0.55,0.3) arc (20:200:0.55 and 0.35);
  \node[main, font=\footnotesize, anchor=south] at (0,0.95) {$\dot\theta_{\mathrm{front}}$ 快 ($I_{\mathrm{front}}$ 小)};
  % rear body: extended (large inertia), slow opposite rotation
  \draw[second, thick, fill=second!10] (2.6,0) ellipse (0.75 and 0.4);
  \node[second, font=\small] at (2.6,0) {后};
  \draw[<-, second, thick] (3.45,-0.3) arc (-20:-200:0.6 and 0.32);
  \node[second, font=\footnotesize, anchor=north] at (2.6,-0.95) {$\dot\theta_{\mathrm{rear}}$ 慢反向 ($I_{\mathrm{rear}}$ 大)};
  % joint
  \draw[thick] (0.45,0) -- (1.85,0);
  \fill (1.15,0) circle (1.6pt);
  \node[font=\footnotesize, anchor=south] at (1.15,0.08) {脊柱关节};
  \node[font=\small, anchor=north] at (1.3,-1.5) {总角动量 $k_G=I_{\mathrm{front}}\dot\theta_{\mathrm{front}}+I_{\mathrm{rear}}\dot\theta_{\mathrm{rear}}=0$，净转角 $\neq0$};
\end{tikzpicture}
\caption{猫翻身：零总角动量下，靠惯量不对称（前蜷缩小、后伸展大）使前后半身反向旋转的转角不抵消，累积出净姿态变化。这是质心模型的盲区——它只看 $k_G=0$，看不见内部角动量重分配。}
\label{fig:cat-flip}
\end{figure}

\begin{insight}[猫翻身是质心模型的盲区]
质心模型只看总角动量 $\boldsymbol{k}_G=0$，无法捕捉内部角动量重分配，于是「预测」猫不能翻身——这与现实矛盾。理解这一点要靠\emph{全身}动力学模型。这正划清了「质心/简化模型」与「全身模型」的能力边界：前者高效但盲于内部自由度，后者完整但昂贵。
\end{insight}

\subsection{工程应用：摆臂与 Flywheel 模型}
\label{subsec:armswing}

行走时腿的摆动产生角动量（左右腿前摆都产生同向角动量），若不补偿躯干会逐渐向后旋转。\textbf{摆臂}：手臂反向摆动产生反向角动量，$k_G^y\approx k_{\mathrm{legs}}^y+k_{\mathrm{arms}}^y=0$。\textbf{Flywheel 模型}：在线性倒立摆（LIPM）基础上增加一个「飞轮」——躯干的角动量自由度，$\dot{k}_G^y=\sum_c(\boldsymbol{p}_c-\boldsymbol{p}_G)\times\boldsymbol{f}_c\cdot\hat{\boldsymbol{y}}$，保留 LIPM 简单性的同时捕捉角动量变化，对快速行走与跑步至关重要。

\begin{pitfall}[概念误区：LIPM 假设 $k_G=0$ 的失效条件]
LIPM 假设角动量恒为零 $k_G=0$，成立条件是慢速行走（$<0.5\,$m/s）、双足支撑时间长、腿摆动角动量小。\emph{严重失效}于快速行走（$>1\,$m/s）、跑步、跳跃、不平地形——此时须升级到质心模型（保留角动量自由度）或 Flywheel 模型。
\end{pitfall}

\begin{pitfall}[思维陷阱：角动量「越小越好」]
在 MPC 中强制 $k_G=0$ 是错的：跳跃前需累积角动量，扰动恢复时适量角动量帮助保持平衡（摆臂效应）。正确做法是\emph{软约束} $\min\|k_G\|^2$，允许必要时偏离零。
\end{pitfall}

\begin{practice}
\begin{enumerate}
  \item 一个四足从 $0.5\,$m 高自由落体，起跳时 $k_G=0$，落地前绕俯仰轴偏转了 $5^\circ$。这是否违反角动量守恒？（提示：考虑空气阻力与内部关节运动。）
  \item 从角动量角度解释为何人形机器人跑步时摆臂幅度大，而四足机器人不需要摆臂。
\end{enumerate}
\end{practice}

% ════════════════════════════════════════════════════════════════
\section{软接触 vs 刚体接触：建模权衡}
\label{sec:contact-soft}

\paragraph{承上启下。}
前面各节默认刚体接触（约束力、精确不穿透），它正是工业控制器与本章 LCP 推导的基础。但仿真器（尤其 MuJoCo、Isaac）常用软接触。为何二者各取所需？本节作一节的对比收束，并为可微接触前沿（\cref{sec:contact-outlook}）铺垫。

\subsection{两类模型}
\label{subsec:two-models}

\textbf{刚体接触}：接触力是\emph{约束力}，不是显式函数，由互补条件隐式确定（即 \cref{sec:contact-complementarity-lcp} 的 LCP）。数学上精确（真不穿透、不滑移），但每一步都是一个优化问题。\textbf{软接触}有几种典型形式。\emph{弹簧--阻尼}：$\lambda_c^n=\max(0,\,k(-\phi)-d\dot\phi)$，外层 $\max(0,\cdot)$ 防止闭合阶段（$\dot\phi<0$）阻尼项使法向力变负（拉力），保证 $\lambda_c^n\ge0$。\emph{Hertz 弹性}：法向弹性力 $\lambda_c^n=k_H|\phi|^{3/2}$ 的 $3/2$ 次律源自 Hertz（1882）的弹性接触理论\cite{hertz1882beruhrung}（球体压入平面的接触力与变形量 $3/2$ 次方成正比）；工程中常叠加一个 Hunt--Crossley（1975）型非线性阻尼项 $d_H|\phi|^{3/2}\dot\phi$ 以耗散碰撞能量\cite{hunt1975coefficient}——\emph{弹性律归 Hertz、阻尼项归 Hunt--Crossley，二者出处不同}。\emph{Drake Hydroelastic}：不是点接触，而是\emph{体积压力场}接触\cite{masterjohn2022velocity}——每个碰撞体被赋予一个压力场（刚度从表面到内部线性增加），两物体压力场在接触区叠加，合力 $\boldsymbol{f}_c=\int_{\Omega}p(\boldsymbol{x})\boldsymbol{n}(\boldsymbol{x})\,dA$ 由积分给出，自然处理面接触、滚动摩擦与 CoP 偏移，代价是需 mesh 离散、参数难校准。

\subsection{对比与工程取舍}
\label{subsec:model-compare}

\begin{table}[H]\centering\small
\caption{刚体接触与软接触的对比}
\label{tab:rigid-vs-soft}
\begin{tabular}{lll}
\toprule
维度 & 刚体接触（LCP）& 软接触（弹簧--阻尼 / Hydroelastic）\\
\midrule
物理精度 & 精确的不穿透 & 允许微小穿透 \\
数学光滑性 & 非光滑（互补约束）& 光滑（连续力函数）\\
计算方法 & LCP/QP 求解 & 常微分积分 \\
时间步 & 可用大步长 & 需小步长（刚性 ODE）\\
参数需求 & 只需 $\mu$ & 需 $k,d$（难校准）\\
面接触 & 需显式建模 & 自然支持（Hydroelastic）\\
控制器适配 & 控制器也用 LCP 思维 & 控制器与仿真不匹配 \\
可微性 & 不可微（需特殊处理）& 天然可微（适合 RL）\\
\bottomrule
\end{tabular}
\end{table}

\noindent \textbf{为何工业控制器用刚体模型？}（1）sim-to-real 一致性——控制器的摩擦锥 QP 与刚体接触模型匹配；（2）参数简单——只需 $\mu$；（3）性能——不需小时间步。\textbf{为何仿真器倾向软接触？}（1）数值稳定——软接触力是连续函数，积分更稳；（2）GPU 并行——软接触不需求解全局 QP，适合大规模并行；（3）RL 训练——RL 需要可微仿真器（策略梯度），软接触天然可微。

\begin{pitfall}[思维陷阱：仿真器选择无所谓？]
不同仿真器的接触模型差异会导致 sim-to-real gap：在某仿真器中训练的策略迁移到另一刚体仿真验证时，行为可能完全不同。软接触的可微性利于 RL 策略梯度，但 sim-to-real、域随机化、仿真器对比等内容属强化学习运控部，本章不展开，详见 RL 运控部相关章。
\end{pitfall}

% ════════════════════════════════════════════════════════════════
\section{本章脉络与展望}
\label{sec:contact-outlook}

\subsection{常见误解汇总}
\label{subsec:misconceptions}

\begin{itemize}
  \item \textbf{外切近似更安全}——错，内切才保守；外切允许的力超出真实锥，可能滑倒（\cref{subsec:cone-error}）。
  \item \textbf{$k=4$ 越多越精确就越好}——$k=4$ 误差约 $27\%$ 与 $\mu$ 测不准（$10\%\sim20\%$）同量级，盲目增 $k$ 是伪精度（\cref{subsec:cone-error}）。
  \item \textbf{互补性就是碰撞检测}——互补是稳态「接触 $\leftrightarrow$ 力」逻辑，碰撞是瞬时事件需冲量模型（\cref{subsec:complementarity}）。
  \item \textbf{半正定 LCP 一定有解}——须附「可行」限定；纯 PSD 本身不无条件保证可解（\cref{thm:lcp-solvable}）。
  \item \textbf{CMM 是 $\mathbf{M}$ 前 6 行}——维度不同、含质心变换，不是截取（\cref{subsec:ccrba}）。
  \item \textbf{CCRBI 是常数}——随构型变（蜷缩 vs 伸展），花滑收臂同理（\cref{subsec:ccrba}）。
  \item \textbf{静平衡可行即动态稳定}——前者是静力学必要条件，后者需控制分析（\cref{subsec:static-equilibrium}）。
  \item \textbf{Hertz 模型含阻尼}——$3/2$ 弹性律归 Hertz，阻尼项归 Hunt--Crossley（\cref{subsec:two-models}）。
\end{itemize}

\begin{table}[H]\centering\small
\caption{本章核心知识点速查}
\label{tab:contact-summary}
\begin{tabular}{lll}
\toprule
知识点 & 核心公式/概念 & 工程应用 \\
\midrule
Signorini 条件 & $\phi\ge0,\ \lambda^n\ge0,\ \phi\lambda^n=0$ & 所有接触约束的基础 \\
库仑摩擦锥 & $\|\boldsymbol{\lambda}^t\|\le\mu\lambda^n$（SOC，凸非线性）& WBC/MPC 核心约束 \\
摩擦锥线性化 & $k$ 边多面体，误差 $\sim1-\cos(\pi/k)$ & 凸 MPC 进 QP \\
$H$/$V$-rep 对偶 & 半空间交 = 生成元非负包 & 检验力 vs 枚举可达力 \\
Wrench cone & 6D 力/力矩耦合约束（$\mathbf{U}\boldsymbol{w}\le0$）& 人形面接触 \\
GIWC 静平衡 & $-\boldsymbol{w}_G^{\mathrm{gravity}}\in$ GIWC & 「能否站稳」的 LP \\
Delassus 算子 & $\mathbf{A}=\mathbf{J}_n\mathbf{M}^{-1}\mathbf{J}_n^\top\succeq0$ & 刚体接触 LCP \\
Fischer--Burmeister & $\sqrt{a^2+b^2}-a-b=0$ & 互补约束软化 \\
CITO & 接触时刻/位置为决策变量 & 离线步态发现 \\
接触 Jacobian 对偶 & $\boldsymbol{\tau}=\mathbf{J}_c^\top\boldsymbol{\lambda}$ & 力从接触到关节空间 \\
CCRBA & $O(N)$ 算 CMM，$\boldsymbol{h}_G=\mathbf{A}_G\dot{\boldsymbol{q}}$ & 质心 MPC/WBC \\
角动量守恒 & 飞行相位 $\dot{\boldsymbol{k}}_G=0$ & 跳跃/翻滚规划 \\
猫翻身 & 零角动量靠惯量不对称 & 全身 vs 质心模型边界 \\
\bottomrule
\end{tabular}
\end{table}

\begin{table}[H]\centering\small
\caption{互补约束三大处理路线对比}
\label{tab:three-routes}
\begin{tabular}{llll}
\toprule
维度 & 预定义模式 & 软化互补 & 接触隐式 TO \\
\midrule
实时性 & 毫秒级 & 秒级 & 分钟级 \\
适应性 & 固定步态 & 需初始猜测 & 自动发现 \\
工业成熟度 & 主流栈 & 研究代码 & 前沿研究 \\
数学难度 & 低 & 中 & 高 \\
代表 & Cheetah、ANYmal & NLP + FB & Posa 2014 \\
\bottomrule
\end{tabular}
\end{table}

\subsection{展望：可微接触前沿}
\label{subsec:differentiable}

传统接触模型（LCP/互补约束）本质上\emph{不可微}——接触状态切换是离散跳变，梯度在切换点不存在或为零，这对端到端学习管线（RL 策略优化、可微仿真反向传播）构成根本障碍。三条可微化路线各有取舍：\emph{penalty / 软接触}（用连续弹簧--阻尼替代刚性，自然可微但时间步受刚度限制）；\emph{隐函数微分}（把 LCP/NCP 的解看作参数的隐函数，用 KKT 条件求导，梯度精确但需解 KKT 系统的逆）；\emph{randomized smoothing}（在接触参数上加噪取期望，得到光滑梯度，无需改仿真器但梯度有偏、方差大）。其核心矛盾是\textbf{物理精度 vs 数学光滑性}：刚性接触更符合现实（地面确实是硬的）但不可微，软接触牺牲精度（人为穿透）换光滑，randomized smoothing 则试图两全。这一前沿的公式推导（如高斯平滑期望与蒙特卡罗梯度估计）属规划部与 RL 运控部，本章不展开，\rebuilt{可微接触前沿（Dojo/ContactNets/randomized smoothing/MJX/SAP 等）联网仍存疑，待核。}

\subsection{启下：本章在全书中的位置}
\label{subsec:next}

本章是腿足/接触子部的理论地基。向后：摩擦锥 QP、GIWC、CMM 全在此打底，下游的腿足凸 MPC 与全身控制（\cref{ch:force-wbc}）、操作空间力控（\cref{ch:operational-space}）在用到时 \cref{ch:contact-mechanics}；接触隐式 TO 的 LCP/Delassus 理论底在 DDP 家族（含 Crocoddyl 接触前向动力学）处复用。横向：LCP 的非对称数值求解、碰撞冲量、约束 ABA 归约束动力学（\cref{ch:constrained-dynamics}），二者以 LCP 标准形与 Delassus 定义为交界，且 CCRBA 与该章的复合刚体算法（CRBA）同源——CCRBA 即 CRBA 的复合惯量递推加质心力变换；空间向量代数的接口在 \cref{ch:spatial-dynamics,ch:geometric-mechanics}。向上：复合系统的统一动力学与多模态 MPC（\cref{ch:unified-multimodal-mpc}）把本章接触机制纳入更大框架。回扣：本章把 \cref{ch:control} 全身控制与凸 MPC 那里一笔带过的接触约束补成了完整理论，二者互为承接。

\paragraph{延伸阅读。}
摩擦锥线性化的工程典范见 Di Carlo 等 2018（MIT Cheetah\,3 凸 MPC）\cite{dicarlo2018cheetah}；CMM/CCRBA 的完整数学见 Orin 等 2013\cite{orin2013centroidal}，其在全身控制中的改进见 Wensing--Orin 2016\cite{wensing2016improved}；LCP 接触时间积分的经典见 Stewart--Trinkle 1996\cite{stewart1996implicit}，理论参考 Cottle--Pang--Stone 1992\cite{cottle1992linear}；wrench cone 与多接触稳定性见 Caron 等 2015 双篇\cite{caron2015stability,caron2015leveraging}；接触隐式 TO 见 Posa 等 2014\cite{posa2014direct} 与 Manchester--Kuindersma 2020\cite{manchester2020variational}；本章主题的腿足优化控制总览见 Wensing 等 2024\cite{wensing2024legged}；接触模型与求解器的系统 benchmark 见 Le Lidec 等 2024\cite{lelidec2024contactbench}。
